\chapter{无穷级数与函数展开}\label{ch:9}
\begin{tcolorbox}[title=章首知识地图,colframe=Ink]
部分和 $\to$ 正项与任意项判敛 $\to$ 幂级数半径、端点与求和 $\to$ 一致收敛 $\to$ Fourier展开与点收敛。
\end{tcolorbox}
\begin{examtip}考研主线：数项判敛、绝对与条件收敛、幂级数收敛域与求和、Fourier系数。常见题型：反例判断、端点检查、逐项运算。竞赛拓展：一致收敛、Dirichlet、积分表示与Fourier选点求和。上述为学习任务分层，具体范围以所报考科目的当年大纲为准。\end{examtip}
\section{级数概念与正项判别}\label{sec:9:1}
\subsection{概念、定理与推导}

\begin{definition}{级数收敛}级数$\sum_{n=1}^\infty a_n$收敛指部分和$S_N=\sum_{n=1}^Na_n$有有限极限。必要条件$a_n\to0$，但不充分。正项级数收敛等价于部分和有界。\end{definition}
\begin{theorem}{比较、比值、根值与积分判别}非负项满足$0\le a_n\le b_n$时，大级数收敛推出小级数收敛。若$a_n\ge0,b_n>0$最终成立，且$\frac{a_n}{b_n}\to c\in(0,\infty)$，则同敛散。比值判别要求$a_n>0$最终成立；比值极限$L=\lim \frac{a_{n+1}}{a_n}$小于1则收敛，大于1则发散；等于1不定。根值上极限$L=\limsup\sqrt[n]{a_n}$同理。若$f$正、递减且连续，则$\sum f(n)$与$\int_1^\infty f(x)\dd x$同敛散。\end{theorem}
\begin{proof}比较来自部分和保序；比值或根值小于1时最终受几何级数控制；比值极限大于1时，项最终以大于1的比例增长；根值上极限大于1时，存在$q>1$及无穷多个$n$使$a_n\ge q^n$。两种情形一般项均不趋零。积分判别用左右矩形夹住积分。\end{proof}
$p$级数$\sum n^{-p}$当且仅当$p>1$收敛。正递减$a_n$的凝聚判别将$\sum a_n$与$\sum2^ka_{2^k}$比较，适合对数因子。

\begin{method}{级数概念与正项判别}条件：先判项是否趋零、是否非负。步骤：指数阶乘用比值，$n$次幂用根值，幂与对数用比较或积分。易错点：比值等于1时不能下结论。\end{method}
\subsection{例题与方法比较}
\begin{example}{基础例题｜调和发散}{难度：$\star$}
\label{ex:9:1}\textbf{题目} 证明$\sum_{n=1}^{\infty}\frac{1}{n}$发散。\par
\textbf{思路} 分组、积分或Cauchy准则。\par
\textbf{解答} 方法一：每组$n=2^{k-1}+1$到$2^k$的和至少$\frac{1}{2}$，部分和无界。方法二：$S_N\ge\int_1^{N+1}\frac{\dd x}{x}=\ln(N+1)\to\infty$。方法三：尾和$\sum_{k=n+1}^{2n}\frac{1}{k}\ge\frac{1}{2}$不满足Cauchy准则。\par
\textbf{方法点评} 分组体现尺度结构，积分给增长下界，Cauchy直接否定收敛。\par
\textbf{易错点或变式} 一般项趋零不能判收敛。
\end{example}
\begin{example}{考研风格例题｜阶乘比值}{难度：$\star\star$}
\label{ex:9:2}\textbf{题目} 判断$\sum \frac{n!}{n^n}$。\par
\textbf{思路} 阶乘使用相邻比。\par
\textbf{解答} $\frac{a_{n+1}}{a_n}=(\frac{n}{n+1})^n\to\e^{-1}<1$，故收敛。另一法：将$\frac{n!}{n^n}=\prod_{k=1}^n(\frac{k}{n})$，前$\lfloor \frac{n}{2}\rfloor$项不超过$\frac{1}{2}$，其余不超过1，故$a_n\le2^{-\lfloor \frac{n}{2}\rfloor}$，比较几何级数收敛。\par
\textbf{方法点评} 直接乘积估计不依赖Stirling公式。\par
\textbf{易错点或变式} 比值中的$n+1$因子会与阶乘变化抵消。
\end{example}
\begin{example}{考研风格例题｜对数临界}{难度：$\star\star$}
\label{ex:9:3}\textbf{题目} 判断$\sum_{n=2}^{\infty}\frac{1}{n(\ln n)^p}$。\par
\textbf{思路} 比值趋1，改用积分。\par
\textbf{解答} 积分令$u=\ln x$，化为$\int_{\ln2}^{\infty}u^{-p}\dd u$，当且仅当$p>1$收敛。也可凝聚得到常数倍$\sum k^{-p}$。\par
\textbf{方法点评} 积分与凝聚都降低一层对数。\par
\textbf{易错点或变式} $p\le0$时项函数最终仍递减；从足够大处开始比较即可。
\end{example}
\section{交错、绝对与条件收敛}\label{sec:9:2}

\begin{method}{命题判断中的四个检验序列}
几何项$2^{-n}$提供快速绝对收敛；$\frac{1}{n}$提供项趋零仍发散；$\frac{(-1)^{n-1}}{n}$提供条件收敛；$(-1)^n$提供分组后为零而原级数发散。有限项增删不影响敛散；两个收敛级数可逐项相加，两个发散级数相加没有统一结论。对任意项级数要避免把正项级数的保序结论直接搬用。
\end{method}

\subsection{概念、定理与推导}

\begin{definition}{绝对与条件收敛}$\sum|a_n|$收敛称绝对收敛；$\sum a_n$收敛而$\sum|a_n|$发散称条件收敛。\end{definition}
\begin{theorem}{Leibniz与Dirichlet}若$b_n\downarrow0$，则$\sum(-1)^{n-1}b_n$收敛，余项绝对值不超过首个被省项$b_{N+1}$。若$\sum_{k=1}^na_k$一致有界且$b_n$单调趋零，则$\sum a_nb_n$收敛。绝对收敛蕴含收敛。\end{theorem}
\begin{proof}交错级数的偶部分和递增、奇部分和递减且相差趋零。Dirichlet用分部求和，把尾和界为常数乘$b_m+\sum_{k\ge m}|b_k-b_{k+1}|$，趋零。绝对收敛由尾和三角不等式与Cauchy准则得到。\end{proof}
Abel判别：若$\sum a_n$收敛、$b_n$单调有界，则$\sum a_nb_n$收敛。条件收敛级数不能任意重排；分组后的收敛也不自动推出原级数收敛。

\begin{method}{交错、绝对与条件收敛}条件：交错判别需要幅度最终单调且趋零。步骤：先查绝对值级数，再用符号抵消；误差由交错余项控制。易错点：仅正负交替并不足够。\end{method}
\subsection{例题与方法比较}
\begin{example}{基础例题｜交错调和}{难度：$\star$}
\label{ex:9:4}\textbf{题目} 判断$\sum\frac{(-1)^{n-1}}{n}$并求和。\par
\textbf{思路} Leibniz判敛，积分或幂级数求值。\par
\textbf{解答} 幅度$\frac{1}{n}\downarrow0$，级数收敛；绝对值调和级数发散，故条件收敛。对有限几何和积分：$\sum_{n=1}^N\frac{(-1)^{n-1}}{n}=\int_0^1\frac{1-(-x)^N}{1+x}\dd x$，误差不超过$\frac{1}{N+1}$，极限为$\ln2$。\par
\textbf{方法点评} 有限和积分避免未经证明的端点逐项操作。\par
\textbf{易错点或变式} 不能用半径内逐项运算直接代端点而不给依据。
\end{example}
\begin{example}{考研风格例题｜反例构造}{难度：$\star\star$}
\label{ex:9:5}\textbf{题目} 若$\sum a_n$收敛，是否必有$\sum a_n^2$收敛？\par
\textbf{思路} 用条件收敛而幅度慢衰减的序列。\par
\textbf{解答} 取$a_n=\frac{(-1)^{n-1}}{\sqrt n}$，原级数由Leibniz收敛，但$\sum a_n^2=\sum\frac{1}{n}$发散。若增加绝对收敛条件，则最终$|a_n|\le1$，有$a_n^2\le|a_n|$而成立。\par
\textbf{方法点评} 反例与修正命题一起记忆，明确缺失的假设。\par
\textbf{易错点或变式} 收敛级数的项可正可负，不能套正项比较。
\end{example}
\begin{example}{考研风格例题｜Dirichlet}{难度：$\star\star$}
\label{ex:9:6}\textbf{题目} 证明当$\theta\notin2\pi\mathbb Z$时$\sum\frac{\sin(n\theta)}{n}$收敛。\par
\textbf{思路} 三角部分和有界。\par
\textbf{解答} 由有限几何和，$|\sum_{k=1}^N\e^{ik\theta}|\le\frac{2}{|1-\e^{i\theta}|}$，其虚部部分和有界；$\frac{1}{n}$单调趋零，Dirichlet成立。\par
\textbf{方法点评} 周期振荡也可产生抵消，不必逐项交错。\par
\textbf{易错点或变式} 这里给的是固定$\theta$的收敛，界在$\theta\to0$时失去统一性。
\end{example}
\section{幂级数、收敛域与求和}\label{sec:9:3}
\subsection{概念、定理与推导}

\begin{theorem}{半径与逐项运算}幂级数$\sum a_n(x-x_0)^n$的半径$R=\frac{1}{\limsup\sqrt[n]{|a_n|}}$，约定$\frac{1}{0}=\infty,\frac{1}{\infty}=0$。在$|x-x_0|<R$绝对收敛，外部发散，端点逐一判定；内部可逐项求导积分，半径保持不变。\end{theorem}
\begin{proof}半径结论来自根值判别。任意闭子区间$|x-x_0|\le r<R$上，可选择$r<\rho<R$建立几何级数控制，从而一致收敛；导数级数同样受带$n$的几何项控制。\end{proof}
基础展开：$\e^x=\sum \frac{x^n}{n!}$；$\sin x=\sum\frac{(-1)^nx^{2n+1}}{(2n+1)!}$；$\cos x=\sum\frac{(-1)^nx^{2n}}{(2n)!}$均全实数成立。$\ln(1+x)=\sum_{n\ge1}\frac{(-1)^{n-1}x^n}{n}$在$-1<x\le1$成立；$(1+x)^\alpha=\sum\binom\alpha n x^n$在$|x|<1$成立。非负整数$\alpha$时后者有限终止。

\begin{method}{幂级数、收敛域与求和}条件：运算先限制在半径内。步骤：先求收敛域，再把$n$与$\frac{1}{n}$对应求导或积分；最后单独处理端点。易错点：求导积分可改变端点收敛。\end{method}
\subsection{例题与方法比较}
\begin{example}{考研风格例题｜半径端点}{难度：$\star\star$}
\label{ex:9:7}\textbf{题目} 求$\sum_{n=1}^{\infty}\frac{x^n}{n}$的收敛域与和。\par
\textbf{思路} 比值判半径，再查$\pm1$。\par
\textbf{解答} $R=1$。$x=1$调和发散，$x=-1$交错收敛，故域$[-1,1)$。内部对几何级数积分得$S(x)=\int_0^x\frac{\dd t}{1-t}=-\ln(1-x)$；端点$x=-1$由有限和积分或Abel端点定理得$-\ln2$。\par
\textbf{方法点评} 收敛域、和函数与端点必须一起给。\par
\textbf{易错点或变式} 不要把$R=1$写成收敛区间$[-1,1]$。
\end{example}
\begin{example}{考研风格例题｜加权几何和}{难度：$\star\star$}
\label{ex:9:8}\textbf{题目} 求$S(x)=\sum_{n=1}^{\infty}nx^n$，$|x|<1$。\par
\textbf{思路} 求导、错位相减或微分方程。\par
\textbf{解答} 方法一：对$\sum_{n\ge0}x^n=\frac{1}{1-x}$求导再乘$x$，$S=\frac{x}{(1-x)^2}$。方法二：绝对收敛允许错位相减，$(1-x)S=\sum_{n\ge1}x^n=\frac{x}{1-x}$。方法三：设$T=\frac{S}{x}$，系数递推给$(1-x)T\prime=2T,T(0)=1$，解得$T=(1-x)^{-2}$。\par
\textbf{方法点评} 求导最短，错位法纯代数，ODE法可推广递推系数。\par
\textbf{易错点或变式} $x=0$时第三法用连续延拓$T(0)=1$。
\end{example}
\begin{example}{考研风格例题｜二阶加权}{难度：$\star\star$}
\label{ex:9:9}\textbf{题目} 求$\sum_{n=1}^{\infty}\frac{n^2}{2^n}$。\par
\textbf{思路} 再次对加权几何和求导。\par
\textbf{解答} $\sum n^2x^n=x\frac\dd{\dd x}\frac{x}{(1-x)^2}=\frac{x(1+x)}{(1-x)^3}$，代$x=\frac{1}{2}$得6。另一法：$n^2=2\binom n2+n$，利用$\sum\binom n2x^n=\frac{x^2}{(1-x)^3}$，加上$\sum nx^n$同得6。\par
\textbf{方法点评} 普通幂权与组合权可相互转换。\par
\textbf{易错点或变式} 求导后必须再乘$x$恢复幂次。
\end{example}
\section{一致收敛与交换运算}\label{sec:9:4}

\begin{center}\begin{tikzpicture}
\begin{axis}[width=.65\textwidth,height=4.3cm,xmin=0,xmax=1,ymin=0,ymax=1,xlabel={$x$},ylabel={$x^n$},legend style={font=\small,at={(.02,.98)},anchor=north west},grid=major,grid style={gray!15}]
\addplot[Ink,domain=0:1,samples=60]{x};\addlegendentry{$n=1$}
\addplot[Teal,domain=0:1,samples=60]{x^5};\addlegendentry{$n=5$}
\addplot[Gold,domain=0:1,samples=100]{x^20};\addlegendentry{$n=20$}
\end{axis}\end{tikzpicture}\par\small 端点附近的误差区不断缩窄，但最大误差没有趋零。
\end{center}

\subsection{概念、定理与推导}

\begin{definition}{一致收敛}函数列$f_n$在$E$上一致收敛到$f$指$\sup_{x\in E}|f_n(x)-f(x)|\to0$；等价于误差门槛中的$N$不依赖$x$。函数项级数的一致收敛指其部分和一致收敛。\end{definition}
\begin{theorem}{M判别与交换定理}若$|u_n(x)|\le M_n$且$\sum M_n<\infty$，则$\sum u_n$一致且绝对收敛。连续函数列一致极限连续；闭区间上可积函数列一致收敛则积分与极限可交换。若$f_n\in C^1[a,b]$，$f_n(x_0)$收敛且$f_n\prime$一致收敛到$g$，则$f_n$一致收敛至$f$且$f\prime=g$。\end{theorem}
\begin{proof}M判别由尾和控制。积分误差不超过区间长度乘上确界误差。导数交换利用$f_n(x)=f_n(x_0)+\int_{x_0}^xf_n\prime$，先交换积分与极限，再用基本定理。\end{proof}
一致收敛保证积分交换，却单独不足以保证求导交换；导数级数需要额外控制。

\begin{method}{一致收敛与交换运算}条件：区别逐点与统一误差。步骤：先求点态极限，再算上确界或找依赖$n$的坏点；正向用M判别。易错点：用固定$x$的估计代替全域一致估计。\end{method}
\subsection{例题与方法比较}
\begin{example}{考研风格例题｜逐点非一致}{难度：$\star\star$}
\label{ex:9:10}\textbf{题目} 研究$x^n$在$[0,1]$与$[0,q]$（$0<q<1$）上的一致收敛。\par
\textbf{思路} 端点产生边界层。\par
\textbf{解答} 在$[0,1]$上，点态极限在$x<1$为0、在1为1，不连续，而各项连续，所以不一致。也可取$x_n=1-\frac{1}{n}$，误差趋$\e^{-1}$。在$[0,q]$上$\sup x^n=q^n\to0$，故一致趋0。\par
\textbf{方法点评} 连续极限判据给快速否定，上确界给定量证明。\par
\textbf{易错点或变式} 开区间$[0,1)$上的上确界仍为1，也不一致。
\end{example}
\begin{example}{考研风格例题｜M判别}{难度：$\star\star$}
\label{ex:9:11}\textbf{题目} 证明$\sum_{n=1}^{\infty}\frac{\sin(nx)}{n^2}$在$\R$上一致收敛。\par
\textbf{思路} 正弦被1统一控制。\par
\textbf{解答} $|\frac{\sin(nx)}{n^2}|\le\frac{1}{n^2}$，控制级数收敛，M判别成立。每项连续，所以和函数连续。导数形式$\sum\frac{\cos(nx)}{n}$不满足同一控制，不能据此宣称可逐项求导。\par
\textbf{方法点评} 判别条件与可得结论要分别列明。\par
\textbf{易错点或变式} 一致收敛不是任意交换导数的许可。
\end{example}
\begin{example}{考研风格例题｜移动尖峰}{难度：$\star\star$}
\label{ex:9:12}\textbf{题目} $f_n(x)=nx(1-x)^n$，$0\le x\le1$，是否一致趋0？\par
\textbf{思路} 极大值位置随$n$移动。\par
\textbf{解答} 每个固定$x$都趋0。求导得$f_n\prime=n(1-x)^{n-1}[1-(n+1)x]$，最大值在$x=\frac{1}{n+1}$，值$\frac n{n+1}(\frac n{n+1})^n\to\e^{-1}$，故不一致。\par
\textbf{方法点评} 坏点依赖$n$是逐点与一致的核心差别。\par
\textbf{易错点或变式} 只有逐点趋零不能控制最大误差。
\end{example}
\section{Fourier级数与综合求和}\label{sec:9:5}

\begin{method}{半区间展开与周期缩放}
在$(0,L)$作奇延拓得到正弦级数，$b_n=(\frac{2}{L})\int_0^Lf(x)\sin(\frac{n\pi x}{L})\dd x$；作偶延拓得到余弦级数，$a_n=(\frac{2}{L})\int_0^Lf(x)\cos(\frac{n\pi x}{L})\dd x$，常数项仍为$\frac{a_0}{2}$。端点收敛值由延拓后的左右极限平均决定。取点求数项和前先确认该点满足收敛定理；想求平方系数和时可考虑Parseval等式。
\end{method}

\subsection{概念、定理与推导}

\begin{definition}{Fourier系数}对$2\pi$周期可积函数，$a_n=\pi^{-1}\int_{-\pi}^{\pi}f(x)\cos nx\dd x$，$b_n=\pi^{-1}\int_{-\pi}^{\pi}f(x)\sin nx\dd x$，级数为$\frac{a_0}{2}+\sum_{n\ge1}(a_n\cos nx+b_n\sin nx)$。\end{definition}
\begin{theorem}{常用点收敛定理，标准结论}若周期函数分段$C^1$且分段数有限，则Fourier级数在每点收敛到$\frac{f(x^-)+f(x^+)}{2}$，端点按周期延拓取左右极限。仅连续不足以保证处处收敛。\end{theorem}
\begin{proof}三角正交性给系数公式。将部分和写成Dirichlet核与$f$的卷积，减去左右极限平均后，分段光滑性使差商可积；振荡积分趋零，得到结论。此证明思路中的振荡积分引理可用分段常函数逼近证明。\end{proof}
半区间$(0,L)$的正弦展开采用奇延拓，余弦展开采用偶延拓，频率为$\frac{n\pi}{L}$。高级几乎处处收敛理论属于进阶内容，不列作考研必备定理。
\begin{extension}在平方可积框架下，Parseval等式为$\pi^{-1}\int_{-\pi}^{\pi}|f|^2=\frac{a_0^2}{2}+\sum_{n\ge1}(a_n^2+b_n^2)$；它依赖三角系统完备性，是进一步学习的标准结论。\end{extension}

\begin{method}{Fourier级数与综合求和}条件：先定周期与奇偶延拓。步骤：算系数，写收敛值，再选合法点求和；不连续点用左右平均。易错点：端点函数值不一定等于级数和。\end{method}
\subsection{例题与方法比较}
\begin{example}{竞赛风格例题｜锯齿展开}{难度：$\star\star\star$}
\label{ex:9:13}\textbf{题目} 求$f(x)=x$在$(-\pi,\pi)$的周期Fourier展开。\par
\textbf{思路} 奇函数仅有正弦项。\par
\textbf{解答} $a_n=0$，分部积分得\[b_n=\dfrac2\pi\int_0^\pi x\sin nx\dd x=\dfrac{2(-1)^{n+1}}n.\]因此$x=2\sum_{n\ge1}\frac{(-1)^{n+1}\sin nx}{n}$（$|x|<\pi$）。在$\pm\pi$处周期左右极限为$\pi,-\pi$，级数和为0。\par
\textbf{方法点评} 系数随$\frac{1}{n}$衰减，与周期延拓的跳跃相一致。\par
\textbf{易错点或变式} 不能在$x=\pi$处把左端仍写成$\pi$。
\end{example}
\begin{example}{竞赛风格例题｜平方展开求和}{难度：$\star\star\star$}
\label{ex:9:14}\textbf{题目} 由$f(x)=x^2$的展开求$\sum_{n=1}^{\infty}\frac{1}{n^2}$。\par
\textbf{思路} 偶函数两次分部，再取端点。\par
\textbf{解答} $a_0=\frac{2\pi^2}{3}$；两次分部得$a_n=\frac{4(-1)^n}{n^2},b_n=0$。故$x^2=\frac{\pi^2}{3}+4\sum\frac{(-1)^n\cos nx}{n^2}$。周期延拓在$\pi$连续，代入得$\pi^2=\frac{\pi^2}{3}+4\sum\frac{1}{n^2}$，所以和为$\frac{\pi^2}{6}$。\par
\textbf{方法点评} 函数展开将数项求和转为选点。\par
\textbf{易错点或变式} 先知道级数绝对收敛，有助于验证端点代入稳定。
\end{example}
\begin{example}{竞赛风格例题｜裂项与积分}{难度：$\star\star\star$}
\label{ex:9:15}\textbf{题目} 求$\sum_{n=1}^{\infty}\frac{1}{n(n+1)}$。\par
\textbf{思路} 有限部分和裂项，或用非负积分。\par
\textbf{解答} 方法一：$\frac{1}{n(n+1)}=\frac{1}{n}-\frac{1}{n+1}$，$S_N=1-\frac{1}{N+1}\to1$。方法二：$\frac{1}{n(n+1)}=\int_0^1x^{n-1}(1-x)\dd x$，非负项交换后得$\int_0^1\frac{1-x}{1-x}\dd x=1$。\par
\textbf{方法点评} 裂项最短，积分表示可推广复杂分母。\par
\textbf{易错点或变式} 交换无穷和与积分需说明非负性或一致收敛等依据。
\end{example}
\begin{warning}项趋零不是级数收敛的充分条件；条件收敛不允许任意重排；幂级数半径不决定端点；连续函数的Fourier级数不保证处处收敛，不能把点收敛、一致收敛与几乎处处收敛混用。\end{warning}


\section{级数的判别、求和与反例体系}\label{mat:series_methods}
\begin{examtip}把数项级数、函数项级数和求和技巧之间的依赖关系补入正文。\end{examtip}
\subsection{Cauchy凝聚与Raabe判别}\label{mat:series_methods:condense}
\begin{method}{Cauchy凝聚与Raabe判别}\textbf{适用条件} 凝聚要求非负递减；Raabe要求正项。\par\textbf{实施与依据} 凝聚比较$\sum a_n$与$\sum2^na_{2^n}$。若$\lim n(\frac{a_n}{a_{n+1}}-1)=r$，则$r>1$收敛、$r<1$发散，$r=1$不定。\end{method}
\begin{example}{专题例题｜Cauchy凝聚与Raabe判别}{难度：$\star\star$}
\textbf{题目} 用凝聚判断$\sum\frac{1}{n\ln n}$，$n\ge2$。\par\textbf{思路与解答} 凝聚项$\frac{2^k}{2^k\ln2^k}=\frac{1}{k\ln2}$，为调和型，故发散。积分判别也给$\ln\ln x$发散。\par\textbf{方法点评与边界} 临界判别无结论时继续比较下一层对数。
\end{example}
\subsection{绝对收敛、重排与Cauchy积}\label{mat:series_methods:rearrange}
\begin{method}{绝对收敛、重排与Cauchy积}\textbf{适用条件} 无穷和次序变化。\par\textbf{实施与依据} 绝对收敛允许任意重排；两绝对收敛级数的Cauchy积绝对收敛且和为乘积。\end{method}
\begin{example}{专题例题｜绝对收敛、重排与Cauchy积}{难度：$\star\star$}
\textbf{题目} 用Cauchy积解释$\e^x\e^y=\e^{x+y}$。\par\textbf{思路与解答} 系数卷积为$\sum_{k=0}^n\frac{x^ky^{n-k}}{k!(n-k)!}=\frac{(x+y)^n}{n!}$，绝对收敛允许重排，得到指数加法公式。\par\textbf{方法点评与边界} 条件收敛时任意重排或乘积可能改变结果甚至发散。
\end{example}
\subsection{加括号与原级数}\label{mat:series_methods:group}
\begin{method}{加括号与原级数}\textbf{适用条件} 从分组和判断原级数。\par\textbf{实施与依据} 原级数收敛可任意保持顺序分组；反向需组内余项统一趋零。\end{method}
\begin{example}{专题例题｜加括号与原级数}{难度：$\star\star$}
\textbf{题目} 为何$1-1+1-1+\cdots$按$(1-1)$分组不能证明原级数和0？\par\textbf{思路与解答} 偶部分和0，奇部分和1，两子列不同。组端点收敛只控制部分子列，组内跳动并未趋零。\par\textbf{方法点评与边界} 判断级数敛散性应考察原部分和，而不能只作形式运算。
\end{example}
\subsection{八种求和入口}\label{mat:series_methods:eight}
\begin{method}{八种求和入口}\textbf{适用条件} 收敛性已确认。\par\textbf{实施与依据} 按裂项、已知和变形、幂级数、逐项求导、逐项积分、微分方程、Fourier与复指数恒等式选择求和入口；每次运算核验条件。\end{method}
\begin{example}{专题例题｜八种求和入口}{难度：$\star\star$}
\textbf{题目} 求$\sum_{n=1}^\infty\frac{1}{n2^n}$。\par\textbf{思路与解答} 方法一：$-\ln(1-x)$在$x=\frac{1}{2}$给$\ln2$。方法二：$\sum x^{n-1}$在$[0,\frac{1}{2}]$一致收敛，逐项积分得$\int_0^{\frac{1}{2}}\frac{\dd t}{1-t}=\ln2$。\par\textbf{方法点评与边界} 不同方法之间会重叠，必须说明核心依据而非仅换名称。
\end{example}
\subsection{系数递推与生成函数ODE}\label{mat:series_methods:odegen}
\begin{method}{系数递推与生成函数ODE}\textbf{适用条件} 系数满足多项式递推。\par\textbf{实施与依据} 把$n a_n$对应$xF\prime$，平移下标时补齐初项。\end{method}
\begin{example}{专题例题｜系数递推与生成函数ODE}{难度：$\star\star$}
\textbf{题目} 由$a_0=1,(n+1)a_{n+1}=a_n$求$F(x)=\sum a_nx^n$。\par\textbf{思路与解答} 递推给$F\prime=F,F(0)=1$，所以$F=\e^x$。也可直接得到$a_n=\frac{1}{n!}$。\par\textbf{方法点评与边界} 形式推导后应核对收敛半径，使生成函数确为解析函数。
\end{example}
\subsection{Abel端点定理}\label{mat:series_methods:abel}
\begin{method}{Abel端点定理}\textbf{适用条件} 幂级数在端点1收敛。\par\textbf{实施与依据} 利用部分和有界及分部求和控制端点尾部。\end{method}
\begin{example}{专题例题｜Abel端点定理}{难度：$\star\star$}
\textbf{题目} 说明$\lim_{r\uparrow1}\sum_{n\ge1}\frac{(-1)^{n-1}r^n}{n}=\ln2$。\par\textbf{思路与解答} 端点交错调和收敛，Abel定理给从内部的极限等于端点和；内部和为$\ln(1+r)$，极限$\ln2$。直接用有限和积分可提供独立证明。\par\textbf{方法点评与边界} 内部函数的有限极限并不自动保证端点级数普通收敛。
\end{example}
\subsection{一致Dirichlet与逐项求导}\label{mat:series_methods:diruniform}
\begin{method}{一致Dirichlet与逐项求导}\textbf{适用条件} 参数化部分和一致有界，衰减因子一致趋零并单调。\par\textbf{实施与依据} 对尾和统一使用分部求和；求导则另需导数级数一致收敛及一点收敛。\end{method}
\begin{example}{专题例题｜一致Dirichlet与逐项求导}{难度：$\star\star$}
\textbf{题目} 证明$\sum\frac{\sin(nx)}{n}$在$[\delta,2\pi-\delta]$一致收敛，$0<\delta<\pi$。\par\textbf{思路与解答} 三角部分和的模被$\frac{2}{|1-\e^{ix}|}$控制，该区间分母有正下界。乘$\frac{1}{n}\downarrow0$，一致Dirichlet成立。\par\textbf{方法点评与边界} 当$\delta\to0$时统一界失效，因此不能直接扩大到含0的整个区间。
\end{example}
\subsection{Weierstrass多项式逼近}\label{mat:series_methods:weierstrass}
\begin{method}{Weierstrass多项式逼近}\textbf{适用条件} 紧区间上的连续函数。\par\textbf{实施与依据} 缩放到$[0,1]$，用Bernstein多项式及模连续性统一控制。\end{method}
\begin{example}{专题例题｜Weierstrass多项式逼近}{难度：$\star\star$}
\textbf{题目} 对$f(x)=x^2$求$B_nf(x)$。\par\textbf{思路与解答} 若$X\sim\operatorname{Bin}(n,x)$，则$B_nf=\mathbb E(\frac{X}{n})^2=x^2+\frac{x(1-x)}{n}$，上确界误差$\frac{1}{4n}$，直接验证一致逼近。\par\textbf{方法点评与边界} 逼近存在性与Taylor级数是否收敛到函数不同；光滑函数未必等于其Taylor级数。
\end{example}
\subsection{调和级数的多种发散机制}\label{mat:series_methods:harmonic}
\begin{method}{调和级数的多种发散机制}\textbf{适用条件} 正项且尺度稳定。\par\textbf{实施与依据} 比较积分、二进分组、尾和Cauchy反例和反证结构。\end{method}
\begin{example}{专题例题｜调和级数的多种发散机制}{难度：$\star\star$}
\textbf{题目} 若假设$H_n\to L$，构造矛盾。\par\textbf{思路与解答} 按偶数项与奇数项拆分，$H_{2n}-H_n=\sum_{k=n+1}^{2n}\frac{1}{k}\ge\frac{1}{2}$；但若$H_n$收敛，该差应趋0，矛盾。这也说明$H_n$不是Cauchy列。\par\textbf{方法点评与边界} 多种证明均说明固定比例区间的尾和无法趋于零。
\end{example}

\section{八种级数求和方法}\label{mat:series_sum8}
\begin{examtip}以下按八种求和方法组织算例，与通用判敛和逐项运算相互配合。\end{examtip}
\subsection{一、利用定义求和}\label{mat:series_sum8:j1}
\begin{method}{一、利用定义求和}\textbf{适用条件} 有限部分和可写成闭式。\par\textbf{实施与依据} 先求$S_N$再令$N\to\infty$，不作未经许可的无穷重排。\end{method}
\begin{example}{专题例题｜一、利用定义求和}{难度：$\star\star$}
\textbf{题目} 求$\sum_{n=1}^\infty\frac{2}{n(n+1)(n+2)}$。\par\textbf{思路与解答} 裂项为$\frac{1}{n(n+1)}-\frac{1}{(n+1)(n+2)}$，部分和$\frac{1}{2}-\frac{1}{(N+1)(N+2)}$，故和$\frac{1}{2}$。\par\textbf{方法点评与边界} 这是按定义取极限，不需要幂级数。
\end{example}
\subsection{二、利用已知级数和公式}\label{mat:series_sum8:j2}
\begin{method}{二、利用已知级数和公式}\textbf{适用条件} 已知收敛级数的删项与子列可合法处理。\par\textbf{实施与依据} 精确对应下标，再减去有限项。\end{method}
\begin{example}{专题例题｜二、利用已知级数和公式}{难度：$\star\star$}
\textbf{题目} 求$\sum_{n=1}^\infty\frac{1}{(n+2)^2}$。\par\textbf{思路与解答} 利用$\sum_{k\ge1}k^{-2}=\frac{\pi^2}{6}$，去掉$k=1,2$，得到$\frac{\pi^2}{6}-\frac{5}{4}$。\par\textbf{方法点评与边界} 下标从$n=1$开始，不能误把$k=2$保留。
\end{example}
\subsection{三、利用幂级数求和}\label{mat:series_sum8:j3}
\begin{method}{三、利用幂级数求和}\textbf{适用条件} 参数点位于收敛半径内部。\par\textbf{实施与依据} 先建立和函数，再代入给定参数。\end{method}
\begin{example}{专题例题｜三、利用幂级数求和}{难度：$\star\star$}
\textbf{题目} 求$\sum_{n=0}^\infty\frac{1}{2^n(n+1)}$。\par\textbf{思路与解答} $\sum \frac{x^n}{n+1}=-\frac{\ln(1-x)}{x}$，在$x=0$连续补1。代$x=\frac{1}{2}$得到$2\ln2$。\par\textbf{方法点评与边界} 把积分生成的$x^{n+1}$除以$x$时要单独解释0点。
\end{example}
\subsection{四、利用Fourier级数求和}\label{mat:series_sum8:j4}
\begin{method}{四、利用Fourier级数求和}\textbf{适用条件} 展开及所选点满足收敛定理。\par\textbf{实施与依据} 先算系数，再根据待求符号选点。\end{method}
\begin{example}{专题例题｜四、利用Fourier级数求和}{难度：$\star\star$}
\textbf{题目} 从$x^2$展开求交错平方倒数和。\par\textbf{思路与解答} $x^2=\frac{\pi^2}{3}+4\sum\frac{(-1)^n\cos nx}{n^2}$，在$x=0$代入得$\sum_{n\ge1}\frac{(-1)^{n-1}}{n^2}=\frac{\pi^2}{12}$。\par\textbf{方法点评与边界} 代$x=\pi$得到非交错和，两点作用不同。
\end{example}
\subsection{五、利用代数运算求和}\label{mat:series_sum8:j5}
\begin{method}{五、利用代数运算求和}\textbf{适用条件} 绝对收敛允许错位相减。\par\textbf{实施与依据} 用$S-xS$降低系数次数，反复操作得到几何级数。\end{method}
\begin{example}{专题例题｜五、利用代数运算求和}{难度：$\star\star$}
\textbf{题目} 求$S=\sum n^2x^n$，$|x|<1$。\par\textbf{思路与解答} $(1-x)S=\sum_{n\ge1}(2n-1)x^n=\frac{2x}{(1-x)^2}-\frac{x}{1-x}$，故$S=\frac{x(1+x)}{(1-x)^3}$。\par\textbf{方法点评与边界} 无穷移位相减先有收敛依据，不能照搬到发散级数。
\end{example}
\subsection{六、利用微分方程求和}\label{mat:series_sum8:j6}
\begin{method}{六、利用微分方程求和}\textbf{适用条件} 系数在求导后形成递推。\par\textbf{实施与依据} 给和函数建立ODE并利用初值唯一性。\end{method}
\begin{example}{专题例题｜六、利用微分方程求和}{难度：$\star\star$}
\textbf{题目} 求$F(x)=\sum_{n=0}^\infty \frac{x^{2n}}{(2n)!}$。\par\textbf{思路与解答} 全实轴绝对收敛且可逐项求导，$F\prime\prime=F,F(0)=1,F\prime(0)=0$，所以$F=\cosh x=\frac{\e^x+\e^{-x}}{2}$。\par\textbf{方法点评与边界} 只解ODE通解而不代初值，会留下不该有的任意常数。
\end{example}
\subsection{七、利用Euler公式求和}\label{mat:series_sum8:j7}
\begin{method}{七、利用Euler公式求和}\textbf{适用条件} $|q|<1$，振荡项带几何权。\par\textbf{实施与依据} 用$\e^{i\theta}$把三角序列变为几何级数，再取实部。\end{method}
\begin{example}{专题例题｜七、利用Euler公式求和}{难度：$\star\star$}
\textbf{题目} 求$\sum_{n=1}^\infty q^n\cos((2n-1)\theta)$。\par\textbf{思路与解答} 复和为$\frac{q\e^{i\theta}}{1-q\e^{2i\theta}}$。乘共轭取实部，结果$\frac{q(1-q)\cos\theta}{1-2q\cos2\theta+q^2}$。\par\textbf{方法点评与边界} 分母是复数模方，始终正；不能只把分子分母分别取实部。
\end{example}
\subsection{八、利用复数求和}\label{mat:series_sum8:j8}
\begin{method}{八、利用复数求和}\textbf{适用条件} $0\le r<1$，复对数可取连续主支。\par\textbf{实施与依据} $\sum \frac{z^n}{n}=-\log(1-z)$，取实部或虚部。\end{method}
\begin{example}{专题例题｜八、利用复数求和}{难度：$\star\star$}
\textbf{题目} 求$\sum_{n=1}^\infty \frac{r^n\sin(n\theta)}{n}$。\par\textbf{思路与解答} 取$z=r\e^{i\theta}$，因$1-r\cos\theta>0$，虚部为$\arctan[\frac{r\sin\theta}{1-r\cos\theta}]$。\par\textbf{方法点评与边界} $r\uparrow1$的端点需要另用Abel定理与点态收敛条件。
\end{example}

\section{Fourier点收敛的理论链与进阶内容}\label{mat:fourier_advanced}
\begin{examtip}本节介绍Fourier收敛理论的主要对象与证明依赖。初等部分给出推导，Carleson--Hunt定理给出结论与证明路线；后者的核心极大估计作为高级标准结论使用。\end{examtip}
\subsection{Dirichlet核与卷积表达}\label{mat:fourier_advanced:kernel}
\begin{method}{Dirichlet核与卷积表达}\textbf{适用条件} $f$为$2\pi$周期可积函数。\par\textbf{实施与依据} 采用$\widehat f(k)=(2\pi)^{-1}\int_{-\pi}^{\pi}f(t)\e^{-ikt}\dd t$；有限和可直接交换积分。\end{method}
\begin{example}{专题例题｜Dirichlet核与卷积表达}{难度：$\star\star$}
\textbf{题目} 推导部分和$S_Nf$的核表达。\par\textbf{思路与解答} $S_Nf(x)=(2\pi)^{-1}\int_{-\pi}^{\pi}D_N(x-t)f(t)\dd t$，其中$D_N(u)=\sum_{k=-N}^N\e^{iku}=\frac{\sin((N+\frac{1}{2})u)}{\sin(\frac{u}{2})}$。$u=0$取连续值$2N+1$。\par\textbf{方法点评与边界} $D_N$有变号，且其$L^1$范数随$N$增长，不能当作统一正平均核。
\end{example}
\subsection{局部收敛与Dini条件}\label{mat:fourier_advanced:localization}
\begin{method}{局部收敛与Dini条件}\textbf{适用条件} 给定点左右极限存在，差商可积。\par\textbf{实施与依据} 将$f(x+t)+f(x-t)-f(x^+)-f(x^-)$除以$t$，使用振荡积分趋零。\end{method}
\begin{example}{专题例题｜局部收敛与Dini条件}{难度：$\star\star$}
\textbf{题目} 若$f$在$x$附近满足Hölder条件$|f(x+t)-f(x)|\le C|t|^\alpha$，$\alpha>0$，说明收敛。\par\textbf{思路与解答} 差商受$C|t|^{\alpha-1}$控制，在0附近可积，Dini型收敛条件成立，故$S_Nf(x)\to f(x)$。\par\textbf{方法点评与边界} 只连续时模连续性除以$t$未必可积，因此不能直接得到每点收敛。
\end{example}
\subsection{Fejér核与Cesàro求和}\label{mat:fourier_advanced:fejer}
\begin{method}{Fejér核与Cesàro求和}\textbf{适用条件} 连续周期函数，或适当可积函数。\par\textbf{实施与依据} $\sigma_Nf=(N+1)^{-1}\sum_{j=0}^NS_jf$，核$F_N(t)=(N+1)^{-1}[\frac{\sin(\frac{(N+1)t}{2})}{\sin(\frac{t}{2})}]^2$。\end{method}
\begin{example}{专题例题｜Fejér核与Cesàro求和}{难度：$\star\star$}
\textbf{题目} 证明连续周期函数的Fejér均值一致趋于自身。\par\textbf{思路与解答} $F_N\ge0$且积分为$2\pi$。误差积分在$|t|<\delta$部分受模连续性控制；其余部分$F_N\le[(N+1)\sin^2(\frac{\delta}{2})]^{-1}$，质量趋零。因此误差一致趋0。\par\textbf{方法点评与边界} 平均部分和收敛不等于原部分和收敛。
\end{example}
\subsection{Hardy--Littlewood极大函数}\label{mat:fourier_advanced:maximal}
\begin{method}{Hardy--Littlewood极大函数}\textbf{适用条件} 局部可积$f$。\par\textbf{实施与依据} $Mf(x)=\sup_{r>0}|B(x,r)|^{-1}\int_{B(x,r)}|f|$；它衡量局部平均而非点值最大。\end{method}
\begin{example}{专题例题｜Hardy--Littlewood极大函数}{难度：$\star\star$}
\textbf{题目} 对$f=\mathbf1_{[0,1]}$解释$Mf$的远端衰减及$L^1$问题。\par\textbf{思路与解答} 当$x>1$，选取包含$[0,1]$且以$x$为中心的区间，其平均不小于常数乘$\frac{1}{x}$。因此$Mf$的尾部不可积；$f\in L^1$并不推出$Mf\in L^1$。\par\textbf{方法点评与边界} 正确端点结论是弱$(1,1)$估计，而不是强$L^1$有界。
\end{example}
\subsection{覆盖引理与弱型估计}\label{mat:fourier_advanced:weak}
\begin{method}{覆盖引理与弱型估计}\textbf{适用条件} $f\in L^1$，阈值$\lambda>0$。\par\textbf{实施与依据} 对$Mf>\lambda$的每点选平均过阈值的球，再选两两不交子族，用有限倍膨胀覆盖原集合。\end{method}
\begin{example}{专题例题｜覆盖引理与弱型估计}{难度：$\star\star$}
\textbf{题目} 解释$|\{Mf>\lambda\}|\le \frac{C_d\|f\|_1}{\lambda}$。\par\textbf{思路与解答} 所选不交球$B_j$满足$\lambda|B_j|<\int_{B_j}|f|$。膨胀覆盖给坏集测度$\le C_d\sum|B_j|\le \frac{C_d\|f\|_1}{\lambda}$。一般情形先取紧子集及有限覆盖再逼近。\par\textbf{方法点评与边界} 相交球不能直接把积分相加而不控制重复次数。
\end{example}
\subsection{Lebesgue微分定理}\label{mat:fourier_advanced:differentiate}
\begin{method}{Lebesgue微分定理}\textbf{适用条件} 局部可积函数，结论在几乎处处成立。\par\textbf{实施与依据} 先对连续函数由一致局部变化证明，再用$L^1$逼近及极大函数弱型估计控制误差。\end{method}
\begin{example}{专题例题｜Lebesgue微分定理}{难度：$\star\star$}
\textbf{题目} 为何可把几乎每点视为局部平均的极限？\par\textbf{思路与解答} 对连续近似$g$，$\lim_{r\to0}\frac1{|B(x,r)|}\int_{B(x,r)}|g(y)-g(x)|=0$。$f-g$的局部平均由$M(f-g)$控制，其超阈值集合测度可由$\|f-g\|_1$任意压小，最终得到$f$的Lebesgue点结论。\par\textbf{方法点评与边界} “几乎处处”不允许改写为“每一点”，修改零测集上的值不会改变积分。
\end{example}
\subsection{分布函数与Marcinkiewicz插值}\label{mat:fourier_advanced:interpolation}
\begin{method}{分布函数与Marcinkiewicz插值}\textbf{适用条件} 次线性算子在两个端点有弱型控制。\par\textbf{实施与依据} 按阈值把$f$拆为大值与小值部分，积分分布函数$\|Tf\|_p^p=p\int_0^\infty\lambda^{p-1}|\{|Tf|>\lambda\}|\dd\lambda$。\end{method}
\begin{example}{专题例题｜分布函数与Marcinkiewicz插值}{难度：$\star\star$}
\textbf{题目} 说明为何$M:L^p\to L^p$，$1<p<\infty$。\par\textbf{思路与解答} 弱$(1,1)$与$L^\infty$范数不增给端点估计。对每个阈值分解$f$后积分可得$\|Mf\|_p\le C_d\frac p{p-1}\|f\|_p$。该常数在$p\downarrow1$时发散，与端点仅弱型相容。\par\textbf{方法点评与边界} 插值的分解依赖阈值，不能只做一次固定分解。
\end{example}
\subsection{Banach连续性原理}\label{mat:fourier_advanced:banach}
\begin{method}{Banach连续性原理}\textbf{适用条件} $T_n:L^p\to L^0$线性且依测度连续，极大算子$T^*f=\sup_n|T_nf|$几乎处处有限。\par\textbf{实施与依据} 由Baire类定理获得极大算子在0处依测度的统一控制，再用稠密子空间推广收敛。\end{method}
\begin{example}{专题例题｜Banach连续性原理}{难度：$\star\star$}
\textbf{题目} 给出从三角多项式推广Fourier收敛的框架。\par\textbf{思路与解答} 三角多项式的部分和最终稳定。对一般$f$取多项式$P$逼近，尾部振幅$\Omega(f)=\limsup_{m,n}|S_mf-S_nf|$满足$\Omega(f)\le2S^*(f-P)$。极大估计使右端大于任意固定阈值的集合测度趋零，故$\Omega(f)=0$几乎处处。\par\textbf{方法点评与边界} 关键缺口是极大估计；仅有稠密性不能自动得到逐点收敛。
\end{example}
\subsection{共轭函数与Hilbert变换}\label{mat:fourier_advanced:hilbert}
\begin{method}{共轭函数与Hilbert变换}\textbf{适用条件} 先对光滑紧支撑函数定义，再按范数延拓。\par\textbf{实施与依据} $Hf(x)=\pi^{-1}\operatorname{p.v.}\int_\R \frac{f(t)}{x-t}\dd t$，对称截去奇点后取极限。周期共轭核对应$\cot(\frac{t}{2})$。\end{method}
\begin{example}{专题例题｜共轭函数与Hilbert变换}{难度：$\star\star$}
\textbf{题目} 说明对称奇性如何消去常数项。\par\textbf{思路与解答} 局部积分改写为$\int_{\varepsilon<|u|<1}\frac{[f(x-u)-f(x)]\dd u}{u}$，常数项因核奇而为零；若$f$局部Lipschitz，分子$O(|u|)$，局部奇性可积。\par\textbf{方法点评与边界} 主值是特定截断极限，不是两个发散单侧积分相加。
\end{example}
\subsection{截断算子与极大Hilbert变换}\label{mat:fourier_advanced:hilbertLp}
\begin{method}{截断算子与极大Hilbert变换}\textbf{适用条件} $H_\varepsilon$为去掉半径$\varepsilon$奇点的算子。\par\textbf{实施与依据} 区分统一$L^p$控制、$H^*f=\sup_\varepsilon|H_\varepsilon f|$控制与主值存在。证明常用Calderón--Zygmund分解及插值。\end{method}
\begin{example}{专题例题｜截断算子与极大Hilbert变换}{难度：$\star\star$}
\textbf{题目} 为何逐个$H_\varepsilon$有界还不直接证明主值存在？\par\textbf{思路与解答} 一致范数界不蕴含点态Cauchy。需要在光滑稠密子类先证明主值收敛，再用极大截断算子估计控制逼近误差，逻辑与Banach连续性原理相同。\par\textbf{方法点评与边界} 不能把$\sup$与积分的次序未经估计直接交换。
\end{example}
\subsection{Carleson极大算子与频率调制}\label{mat:fourier_advanced:carleson}
\begin{method}{Carleson极大算子与频率调制}\textbf{适用条件} 振荡奇异积分，频率参数变化。\par\textbf{实施与依据} 模型$Cf(x)=\sup_{\xi\in\R}|\operatorname{p.v.}\int \frac{f(t)\e^{i\xi(x-t)}}{x-t}\dd t|$；此处$\xi$是局部角频率参数；用全书实轴变换的频率记号时写作$\xi=2\pi\eta$，上确界不变。有限区间版本需另外明确核的归一化。\end{method}
\begin{example}{专题例题｜Carleson极大算子与频率调制}{难度：$\star\star$}
\textbf{题目} 解释它为何比单个Hilbert变换困难。\par\textbf{思路与解答} 固定$\xi$时可由调制与Hilbert变换联系，调制保持$L^p$范数。但取所有频率的上确界不能由单个频率的范数界直接推出，需要时频尺度选择和极大不等式。\par\textbf{方法点评与边界} 频率与位置共同变化正是点收敛问题的核心。
\end{example}
\subsection{基本分解：近场、平均与抵消}\label{mat:fourier_advanced:basicstep}
\begin{method}{基本分解：近场、平均与抵消}\textbf{适用条件} 振荡奇异积分的区间分割。\par\textbf{实施与依据} 把振荡奇异积分分为含$x$的近邻积分、远区块平均和均值为零的余项；对远余项利用核的光滑性。\end{method}
\begin{example}{专题例题｜基本分解：近场、平均与抵消}{难度：$\star\star$}
\textbf{题目} 若$b_J$支撑在区间$J$且$\int_Jb_J=0$，估计远处积分。\par\textbf{思路与解答} 取区间中心$c_J$，有$\int_J\frac{b_J(t)}{x-t}\dd t=\int_Jb_J(t)[\frac{1}{x-t}-\frac{1}{x-c_J}]\dd t$。当$|x-c_J|\ge2|J|$，绝对值$\le \frac{C|J|\|b_J\|_1}{|x-c_J|^2}$。\par\textbf{方法点评与边界} 减去常数核利用了零平均，这是比粗估$\frac{1}{\text{距离}}$更强的衰减。
\end{example}
\subsection{二进分割、局部范数与坏集}\label{mat:fourier_advanced:dyadic}
\begin{method}{二进分割、局部范数与坏集}\textbf{适用条件} 多尺度局部平均需要组织。\par\textbf{实施与依据} 将区间不断二等分；按平均超阈值选择极大坏区间，极大性给不交性。\end{method}
\begin{example}{专题例题｜二进分割、局部范数与坏集}{难度：$\star\star$}
\textbf{题目} 给$f\in L^1(I)$、$\lambda>|I|^{-1}\int_I|f|$，估计坏区间总长度。\par\textbf{思路与解答} 每个极大二进坏区间$J$有$\lambda|J|<\int_J|f|$，两两不交，所以$\sum|J|\le\frac{\|f\|_1}{\lambda}$；好区外二进平均不超过阈值，Lebesgue微分将控制传到几乎每点。\par\textbf{方法点评与边界} 局部范数、区间分割与基本步均围绕可累加的局部控制，不只是记一个核公式。
\end{example}
\subsection{频率选择、允许配对与例外集}\label{mat:fourier_advanced:selection}
\begin{method}{频率选择、允许配对与例外集}\textbf{适用条件} Carleson--Hunt定理的多尺度证明。\par\textbf{实施与依据} 限制型弱估计先对$|f|=\mathbf1_A$控制；通过坏集剔除、尺度长度估计、频率平移及允许配对选择，累加基本步，再用插值与稠密推广。\end{method}
\begin{example}{专题例题｜频率选择、允许配对与例外集}{难度：$\star\star$}
\textbf{题目} 说明“先控制例外集”与“每点统一估计”的区别。\par\textbf{思路与解答} 若对每个阈值$\lambda$可找到$|E_\lambda|\le C\lambda^{-p}\|f\|_p^p$且在其外$S^*f\le\lambda$，则这是分布意义的弱型估计，而非在全空间给同一个逐点上界。实现这类控制通常需要频率选择、区间配对和多尺度分解。\par\textbf{方法点评与边界} 这些组合构造须另作论证；本卡只概述证明路线，深层弱型界作为高级标准结论使用。
\end{example}
\begin{theorem}{Carleson--Hunt点收敛：进阶标准结论}在圆周上，$1<p<\infty$时Fourier部分和极大算子满足$\|S^*f\|_p\le C_p\|f\|_p$，因此每个$f\in L^p$的Fourier级数几乎处处收敛到$f$。\end{theorem}
\begin{proof}证明路线：先建立振荡奇异积分的限制型弱估计，借多尺度分解与频率选择控制其上确界；插值获得强$L^p$估计；三角多项式的部分和从某阶起等于原多项式，稠密逼近与极大控制推广到$L^p$；再用$L^p$范数收敛识别点态极限为$f$。前述限制型估计是主证明的核心；这里将该估计作为标准结论，再给出其后续推论。\end{proof}
\subsection{三种收敛层次的边界}\label{mat:fourier_advanced:convergence}
\begin{method}{三种收敛层次的边界}\textbf{适用条件} Fourier范数、一致与几乎处处结论。\par\textbf{实施与依据} 分别检查函数正则性、所处空间和所用求和方式。\end{method}
\begin{example}{专题例题｜三种收敛层次的边界}{难度：$\star\star$}
\textbf{题目} 连续函数的Fourier级数是否必然处处收敛？\par\textbf{思路与解答} 不必然；存在连续函数的Fourier级数在某点发散。连续函数属于任意有限$p$的$L^p$，Carleson--Hunt保证几乎处处收敛，但允许零测异常集；Fejér均值却一致收敛。三结论并不矛盾。\par\textbf{方法点评与边界} 不能把“几乎处处”或“Cesàro求和”省略掉。
\end{example}
\subsection{分布函数、递减重排与Lorentz空间}\label{mat:fourier_advanced:rearrangement}
\begin{method}{分布函数、递减重排与Lorentz空间}\textbf{适用条件} 采用总测度1的空间。\par\textbf{实施与依据} $d_f(\lambda)=\mu\{|f|>\lambda\}$，$f^*(t)=\inf\{\lambda:d_f(\lambda)\le t\}$，$f^{**}(t)=t^{-1}\int_0^tf^*(s)\dd s$。Lorentz范数可按$\int \frac{t^{\frac{1}{p}}f^*(t)\dd t}{t}$或其上确界定义。\end{method}
\begin{example}{专题例题｜分布函数、递减重排与Lorentz空间}{难度：$\star\star$}
\textbf{题目} 计算指示函数$\mathbf1_A$的重排。\par\textbf{思路与解答} 若$\mu(A)=a$，则$f^*(t)=1$（$0<t<a$），其余为0；故弱$L^p$准范数$\sup t^{\frac{1}{p}}f^*(t)=a^{\frac{1}{p}}$，$L^{p,1}$积分型准范数为$p a^{\frac{1}{p}}$，归一化可除$p$。\par\textbf{方法点评与边界} 不同教材常数归一化不同，本册显式写出采用的表达式。
\end{example}
\subsection{接近L1的空间与Antonov路线}\label{mat:fourier_advanced:endpoint}
\begin{method}{接近L1的空间与Antonov路线}\textbf{适用条件} 考察$L^1$端点附近的额外可积性。\par\textbf{实施与依据} 采用$\log(e+t)$等正则化对数定义Orlicz型条件；区分强、弱型与几乎处处收敛。\end{method}
\begin{example}{专题例题｜接近L1的空间与Antonov路线}{难度：$\star\star$}
\textbf{题目} 解释为何引入$L\log L$及更细的对数空间。\par\textbf{思路与解答} 普通$L^1$存在Fourier发散反例，而$p>1$的强型界接近1时常数变坏。分布函数、原子分解和近$L^1$条件把这种常数增长转为端点控制。进一步的端点理论涉及$Q,Q_A$与Antonov的$L\log L\log\log\log L$型充分条件。\par\textbf{方法点评与边界} 这里没有把端点充分条件误写成必要充分条件，也没有声称所有$L^1$函数都收敛。
\end{example}
\subsection{有限Fourier系数的集合逼近}\label{mat:fourier_advanced:antonov}
\begin{method}{有限Fourier系数的集合逼近}\textbf{适用条件} $0\le f\le a\mathbf1_A$，只要求有限个部分和精度。\par\textbf{实施与依据} 把圆周分成足够细区间，每块选$B_k\subset A$使$a|B_k|=\int_{A\cap J_k}f$，利用有限个连续核的一致连续性。\end{method}
\begin{example}{专题例题｜有限Fourier系数的集合逼近}{难度：$\star\star$}
\textbf{题目} 为什么可用$a\mathbf1_B$逼近有限个$S_jf$？\par\textbf{思路与解答} 每块质量完全匹配，核乘常数的主项抵消；误差不超过各核在小块上的振幅乘总质量。对$j\le N$只有有限个核，统一选网格即可使误差小于$\varepsilon$。这给出有限频率条件下的集合替换机制。\par\textbf{方法点评与边界} 这是有限频率逼近，不等于$f$与$a\mathbf1_B$在所有点相近。
\end{example}

\section{端点函数空间与弱型求和细节}\label{mat:endpoint_spaces}
\begin{examtip}本节给出Q、QA等空间的定义，并说明它们在端点估计中的作用。\end{examtip}
\subsection{原子上包络空间Q}\label{mat:endpoint_spaces:qspace}
\begin{method}{原子上包络空间Q}\textbf{适用条件} 圆周归一化测度$\mu(X)=1$。\par\textbf{实施与依据} 若$|f|\le\sum_{j\ge1}a_j\mathbf1_{A_j}$，$a_j>0,m_j=\mu(A_j)$，以$\sum a_jm_j\log(\frac{e}{m_j})(1+\log j)$的下确界定义$\|f\|_Q$。\end{method}
\begin{example}{专题例题｜原子上包络空间Q}{难度：$\star\star$}
\textbf{题目} 估计$f=a\mathbf1_A$的Q准范数。\par\textbf{思路与解答} 只取$j=1$的一项分解，得到$\|f\|_Q\le a\mu(A)\log(\frac{e}{\mu(A)})$。限制型弱估计对单原子给相同规模，弱$L^1$的对数凸求和再控制一般分解。\par\textbf{方法点评与边界} 这给上界，不声称单项分解必定达到定义中的下确界。
\end{example}
\subsection{Antonov空间QA}\label{mat:endpoint_spaces:qaspace}
\begin{method}{Antonov空间QA}\textbf{适用条件} 同一概率测度空间，采用非负有界分解。\par\textbf{实施与依据} 以$|f|\le\sum f_j,f_j\ge0$的所有分解上$\sum(1+\log j)\|f_j\|_1\log(\frac{e\|f_j\|_\infty}{\|f_j\|_1})$的下确界定义$\|f\|_{Q_A}$；零项按0处理。\end{method}
\begin{example}{专题例题｜Antonov空间QA}{难度：$\star\star$}
\textbf{题目} 给出单个有界函数的QA上界。\par\textbf{思路与解答} 取唯一分量$f_1=|f|$，得到$\|f\|_{Q_A}\le\|f\|_1\log(\frac{e\|f\|_\infty}{\|f\|_1})$。利用有限Fourier系数的集合逼近把每分量替成等质量原子，再传递Q空间的弱型估计。\par\textbf{方法点评与边界} 端点弱型估计为$\mu\{S^*f>y\}\le \frac{C\|f\|_{Q_A}}{y}$；定义有限性与这个算子结论是不同层次。
\end{example}
\subsection{弱L1的对数加权求和}\label{mat:endpoint_spaces:weakseries}
\begin{method}{弱L1的对数加权求和}\textbf{适用条件} 弱$L^1$准范数不满足普通三角不等式。\par\textbf{实施与依据} 对可数和使用$\|\sum f_j\|_{1,\infty}\le C\sum(1+\log j)\|f_j\|_{1,\infty}$等对数凸控制。\end{method}
\begin{example}{专题例题｜弱L1的对数加权求和}{难度：$\star\star$}
\textbf{题目} 若$\|f_j\|_{1,\infty}\le2^{-j}$，说明右端有限。\par\textbf{思路与解答} $\sum(1+\log j)2^{-j}$收敛，因此上述定理保证级数和具有弱$L^1$控制。普通$\sum\|f_j\|_{1,\infty}<\infty$本身不能未经论证就按Banach范数三角不等式使用。\par\textbf{方法点评与边界} 对数权重来自坏集测度分配与熵型估计，是Q与QA定义的一部分。
\end{example}

\section{Fourier变换、反演、卷积与应用}\label{mat:fourier_transform}
\begin{examtip}统一周期级数与实轴积分，固定$\widehat f(\xi)=\int f(x)\e^{-2\pi i x\xi}\dd x$的归一化。\end{examtip}
\subsection{Gaussian变换}\label{mat:fourier_transform:gaussian}
\begin{method}{Gaussian变换}\textbf{适用条件} $a>0$，积分绝对收敛。\par\textbf{实施与依据} 对频率求导、分部积分得到一阶ODE，再用零频Gaussian积分定常数。\end{method}
\begin{example}{专题例题｜Gaussian变换}{难度：$\star\star$}
\textbf{题目} 求$f(x)=\e^{-\pi a x^2}$的变换。\par\textbf{思路与解答} $\widehat f\prime(\xi)=-\frac{2\pi\xi\widehat f(\xi)}{a}$，且$\widehat f(0)=a^{-\frac{1}{2}}$，故$\widehat f(\xi)=a^{-\frac{1}{2}}\e^{-\frac{\pi\xi^2}{a}}$。\par\textbf{方法点评与边界} 归一化不同会改变$2\pi$与平方指数系数。
\end{example}
\subsection{卷积、平移与微分}\label{mat:fourier_transform:convolution}
\begin{method}{卷积、平移与微分}\textbf{适用条件} 先在Schwartz函数或足够可积光滑类上。\par\textbf{实施与依据} $\widehat{f*g}=\widehat f\,\widehat g$；$\widehat{f\prime}=2\pi i\xi\widehat f$；平移产生相位，调制产生频移。\end{method}
\begin{example}{专题例题｜卷积、平移与微分}{难度：$\star\star$}
\textbf{题目} 证明卷积变换公式。\par\textbf{思路与解答} 由绝对可积换序，令$u=x-y$，$\int\int f(x-y)g(y)\e^{-2\pi ix\xi}\dd y\dd x$分离为两个变换之积。\par\textbf{方法点评与边界} 交换积分须可积；对一般分布应先说明所在空间。
\end{example}
\subsection{反演与Plancherel}\label{mat:fourier_transform:plancherel}
\begin{method}{反演与Plancherel}\textbf{适用条件} 先对Schwartz函数，随后按$L^2$密度延拓。\par\textbf{实施与依据} 反演$f(x)=\int\widehat f(\xi)\e^{2\pi ix\xi}\dd\xi$；$\|f\|_2=\|\widehat f\|_2$。\end{method}
\begin{example}{专题例题｜反演与Plancherel}{难度：$\star\star$}
\textbf{题目} 解释为何一般$L^2$变换不一定由绝对收敛积分定义。\par\textbf{思路与解答} $L^2$不蕴含$L^1$，所以积分定义可能不绝对收敛。选Schwartz函数列在$L^2$逼近，由等距性使其变换为Cauchy列，极限定义$L^2$Fourier变换，与原积分定义在交集上相容。\par\textbf{方法点评与边界} 范数意义的反演与每点绝对积分反演不是同一结论。
\end{example}
\subsection{热方程与Gaussian核}\label{mat:fourier_transform:heat}
\begin{method}{热方程与Gaussian核}\textbf{适用条件} 初值足够可积，$t>0$。\par\textbf{实施与依据} 空间变换把$u_t=u_{xx}$化为$\partial_t\widehat u=-4\pi^2\xi^2\widehat u$。\end{method}
\begin{example}{专题例题｜热方程与Gaussian核}{难度：$\star\star$}
\textbf{题目} 写出全轴热方程解。\par\textbf{思路与解答} $\widehat u(\xi,t)=\e^{-4\pi^2t\xi^2}\widehat u_0(\xi)$，逆变换为$u(x,t)=(4\pi t)^{-\frac{1}{2}}\int\e^{-\frac{(x-y)^2}{4t}}u_0(y)\dd y$。核质量1，$t\downarrow0$构成近似恒等。\par\textbf{方法点评与边界} 初值的恢复按连续点、$L^p$或几乎处处等适当意义陈述。
\end{example}
\subsection{半轴正弦、余弦变换}\label{mat:fourier_transform:sinecosine}
\begin{method}{半轴正弦、余弦变换}\textbf{适用条件} 边界$x=0$问题。\par\textbf{实施与依据} 奇延拓对应正弦变换，偶延拓对应余弦变换；系数按固定归一化记录。\end{method}
\begin{example}{专题例题｜半轴正弦、余弦变换}{难度：$\star\star$}
\textbf{题目} 计算$\int_0^\infty\e^{-ax}\cos bx\dd x$及正弦版，$a>0$。\par\textbf{思路与解答} 取$\int_0^\infty\e^{-(a-ib)x}\dd x=\frac{1}{a-ib}$的实虚部，分别得到$\frac{a}{a^2+b^2}$与$\frac{b}{a^2+b^2}$。\par\textbf{方法点评与边界} 指数衰减保证边界项消失及绝对可积。
\end{example}
\subsection{实轴Fourier积分的极大截断}\label{mat:fourier_transform:realmax}
\begin{method}{实轴Fourier积分的极大截断}\textbf{适用条件} $1<p<\infty$，实轴上的对称频率截断。\par\textbf{实施与依据} 对对称频率截断的反演积分取上确界，由调制奇异积分控制，并用极大估计与稠密逼近得到收敛。\end{method}
\begin{example}{专题例题｜实轴Fourier积分的极大截断}{难度：$\star\star$}
\textbf{题目} 为何有限频带积分的收敛也需要极大截断估计？\par\textbf{思路与解答} 逐个截断的范数控制不足以保证频带趋无穷时逐点Cauchy。对光滑函数先有反演，再用统一极大截断控制坏集，才能推广到较大$L^p$类。该逻辑与Fourier级数相同。\par\textbf{方法点评与边界} 实轴频率连续，圆周频率离散；转移时不能直接把求和号换成积分号而省去归一化。
\end{example}




\section{判敛逻辑、调和级数多证与求和深化}\label{deep:series_source_deep}

\begin{examtip}同一数学机制的证明合并比较，保留反例、条件与证明差异。以Riemann假设为前提的条件推理不能当作无条件初等证明。\end{examtip}

\subsection{必要条件、括号、乘积与反例}\label{deep:series_source_deep:logic}

\begin{method}{必要条件、括号、乘积与反例}\textbf{适用条件} 先区分数列$a_n$与级数$\sum a_n$，再区分正项与任意项。\par\textbf{步骤与依据} 否定命题时优先尝试调和、交错、稀疏项和配对抵消；对乘积先看一列是否有界。\par\textbf{易错点} 从正项结论删掉正性，往往立刻失效。\end{method}

\begin{example}{考研风格｜括号不可逆}{难度：$\star\star\star$}\label{deep:ex:series_source_deep:block}\textbf{题目} 说明加括号后的级数收敛不能推出原级数收敛，并给充分的补充条件。\par\textbf{思路} 收敛子列的部分和不能控制括号内部的振幅。\par\textbf{详细解答} $1-1+1-1+\cdots$按$(1-1)+(1-1)+\cdots$分组后和为0，原部分和在0与1间跳动。若分组长度统一有界且$a_n\to0$，则从任意部分和到最近分组端点只差有限个尾项，其绝对值$\le K\max_{j\ge n-K}|a_j|\to0$，于是分组和收敛能推出原级数收敛。\par\textbf{方法点评} 反例说明一般不成立，补充条件说明失效发生在何处。\par\textbf{变式与易错点} 若分组长度无界，即使$a_n\to0$也可能在组内累积很大振幅。\end{example}

\begin{example}{考研风格｜逐项乘积与Cauchy乘积}{难度：$\star\star\star$}\label{deep:ex:series_source_deep:product}\textbf{题目} 设$\sum a_n$收敛，$\sum|b_n|$收敛，证明$\sum|a_nb_n|$收敛，并辨析两条件收敛级数的逐项乘积。\par\textbf{思路} 收敛级数的一般项必有界。\par\textbf{详细解答} $a_n\to0$，所以存在$M$使$|a_n|\le M$，从而$\sum|a_nb_n|\le M\sum|b_n|<\infty$。但取$a_n=b_n=\frac{(-1)^{n-1}}{\sqrt n}$，两个级数均由Leibniz收敛，乘积为$\frac{1}{n}$而发散。这里的“逐项乘积”不同于Cauchy乘积$c_n=\sum_{k=0}^na_kb_{n-k}$，不能混用定理。\par\textbf{方法点评} 条件增补与反例构造需要相互配合。\par\textbf{变式与易错点} 绝对收敛性控制的是整个绝对值和，不是只要求$b_n\to0$。\end{example}

\begin{example}{竞赛风格｜收敛不推出$na_n\to0$}{难度：$\star\star\star$}\label{deep:ex:series_source_deep:sparse}\textbf{题目} 构造正项收敛级数，使$na_n$不趋0；再说明单调性如何改变结论。\par\textbf{思路} 将大项放在稀疏指标。\par\textbf{详细解答} 取$a_n=\frac{1}{n^2}$，但当$n=2^k$时改为$a_{2^k}=\frac{1}{2^k}$。总和$\le\sum n^{-2}+\sum2^{-k}<\infty$，而$2^ka_{2^k}=1$。若另有$a_n\downarrow0$，则$\frac n2a_n\le\sum_{k=\lfloor \frac{n}{2}\rfloor}^na_k\to0$，所以$na_n\to0$。\par\textbf{方法点评} 稀疏构造解释正项级数和单调正项级数的差异。\par\textbf{变式与易错点} 比值、根值判别的反命题也可用有稀疏波动的项否定。\end{example}

\subsection{凝聚、Raabe、Dirichlet与Abel}\label{deep:series_source_deep:tests}

\begin{method}{凝聚、Raabe、Dirichlet与Abel}\textbf{适用条件} 正项缓慢变化时考虑凝聚与比值精化；振荡项考虑部分和有界性。\par\textbf{步骤与依据} 先用粗判别，临界值1时展开下一阶；交错结构用分部求和。\par\textbf{易错点} 比值等于1只表示该判别失效，并非级数一定发散。\end{method}

\begin{theorem}{Cauchy凝聚判别}若$a_n\ge0$且递减，则$\sum a_n$与$\sum_{k\ge0}2^ka_{2^k}$同敛散。\end{theorem}\begin{proof}在块$2^k\le n<2^{k+1}$，有$2^ka_{2^{k+1}}\le\sum a_n\le2^ka_{2^k}$。各块相加，左界是凝聚和移位后的一半，右界是凝聚和，故同敛散。\end{proof}

\begin{example}{考研风格｜对数级数的两种判别}{难度：$\star\star\star$}\label{deep:ex:series_source_deep:logseries}\textbf{题目} 判定$\sum_{n\ge2}\frac{1}{n(\ln n)^p}$。\par\textbf{思路} 积分判别或凝聚将对数变成幂。\par\textbf{详细解答} 积分法令$t=\ln x$，$\int_2^\infty\frac{\dd x}{x(\ln x)^p}=\int_{\ln2}^\infty t^{-p}\dd t$，当且仅当$p>1$收敛。凝聚法在项最终递减时，$2^ka_{2^k}=\frac{1}{(k\ln2)^p}$，同样得到$p>1$；若$p\le0$也可直接与调和级数比较发散。\par\textbf{方法点评} 两法揭示相同的尺度变换。\par\textbf{变式与易错点} 从$n=1$开始会遇到$\ln1=0$，必须调整首项。\end{example}

\begin{theorem}{分部求和与Dirichlet判别}若$A_n=\sum_{k=1}^na_k$有界，$b_n$单调趋0，则$\sum a_nb_n$收敛。\end{theorem}\begin{proof}有限分部求和为$\sum_{k=m}^na_kb_k=(A_n-A_{m-1})b_n+\sum_{k=m}^{n-1}(A_k-A_{m-1})(b_k-b_{k+1})$。若$|A_k|\le M$且$b_k\downarrow0$，尾和绝对值$\le2M b_n+2M(b_m-b_n)\le2Mb_m\to0$，由Cauchy准则得收敛。其他单调趋零情形以符号调整。\end{proof}

\begin{example}{竞赛风格｜振荡级数与端点}{难度：$\star\star\star$}\label{deep:ex:series_source_deep:oscillation}\textbf{题目} 判定$\sum_{n\ge1}\frac{\sin(n\theta)}{n}$，并求$0<\theta<2\pi$时的和。\par\textbf{思路} 先用几何和证明部分和有界，再用Abel参数求值。\par\textbf{详细解答} $\sum_{k=1}^n\e^{ik\theta}=\frac{\e^{i\theta}(1-\e^{in\theta})}{1-\e^{i\theta}}$，当$\theta$不是$2\pi$倍数时部分和有界，Dirichlet判别给收敛。对$0<r<1$，$\sum \frac{r^n\e^{in\theta}}{n}=-\log(1-r\e^{i\theta})$。令$r\uparrow1$，Abel定理识别普通和，虚部为$\frac{\pi-\theta}{2}$。$\theta=0,2\pi$时每项为0，和为0，不能代入区间内公式的单侧值。\par\textbf{方法点评} 复数求和给出表达式，端点值还须由收敛定理确认。\par\textbf{变式与易错点} 复对数的辐角必须连续选取；不能把$r=1$直接代进仅对$|z|<1$证明的级数。\end{example}

\subsection{调和级数的证明族：分块、积分与均值}\label{deep:series_source_deep:harmonic}

\begin{method}{调和级数的证明族：分块、积分与均值}\textbf{适用条件} 对正项和可比较有限部分和，并在最后取极限。\par\textbf{步骤与依据} 构造每块统一正下界；或者构造尾和不趋0；避免对发散无穷和直接做代数。\par\textbf{易错点} 写$S>S$须先假设$S$有限并证明严格差保留。\end{method}

\begin{example}{基础到竞赛｜调和级数四证}{难度：$\star\star$}\label{deep:ex:series_source_deep:harmonic4}\textbf{题目} 给出调和级数发散的四种证明。\par\textbf{思路} 采用分块、积分、Cauchy尾和和Cesàro均值。\par\textbf{详细解答} \textbf{分块：}$\sum_{k=2^{j-1}+1}^{2^j}\frac{1}{k}\ge\frac{1}{2}$，所以$H_{2^m}\ge1+\frac{m}{2}$。\par\textbf{积分：}$H_n\ge\int_1^{n+1}\frac{\dd x}{x}=\ln(n+1)\to\infty$。\par\textbf{尾和：}$H_{2n}-H_n\ge\frac{1}{2}$，不满足Cauchy准则。\par\textbf{均值反证：}若$H_n\to S$，则其Cesàro均值也趋$S$；但$n^{-1}\sum_{k=1}^nH_k=(1+\frac{1}{n})H_n-1\to S-1$，矛盾。\par\textbf{方法点评} 分块、积分比较、Cauchy尾和与均值反证给出四种互补的证明思路。\par\textbf{变式与易错点} 最后一个恒等式应先交换有限三角形求和，不涉及发散级数换序。\end{example}

\begin{example}{竞赛风格｜任意底数与Fibonacci分块}{难度：$\star\star\star$}\label{deep:ex:series_source_deep:harmonicblocks}\textbf{题目} 分别按$m^j$和Fibonacci指标分块证明调和级数发散。\par\textbf{思路} 块长除以块末指标给统一下界。\par\textbf{详细解答} 固定整数$m\ge2$，块$(m^{j-1},m^j]$至少贡献$\frac{m^j-m^{j-1}}{m^j}=1-\frac{1}{m}$。对$F_{j+1}=F_j+F_{j-1}$，块$(F_j,F_{j+1}]$贡献至少$\frac{F_{j-1}}{F_{j+1}}\to\phi^{-2}>0$，故两种分块均发散。\par\textbf{方法点评} 两种分块具有同一机制：指数增长块长与终点同阶。\par\textbf{变式与易错点} 分块必须不遗漏尾部且每个整数只计一次。\end{example}

\begin{example}{竞赛风格｜尾和的精确极限与一般分块}{难度：$\star\star\star$}\label{deep:ex:series_source_deep:harmonictail}\textbf{题目} 证明$H_{2n}-H_n\to\ln2$；并解释为何非减块长也能反证调和收敛。\par\textbf{思路} 先积分夹逼，再用块长的单调性。\par\textbf{详细解答} $\int_{n+1}^{2n+1}\frac{\dd x}{x}\le H_{2n}-H_n\le\int_n^{2n}\frac{\dd x}{x}$，两界趋$\ln2$。设正整数块长$i_j$非减且$i_1\ge2$、$I_j=\sum_{\ell\le j}i_\ell$，则$I_j\le ji_j$，第$j$块$\sum_{n=I_{j-1}+1}^{I_j}\frac{1}{n}\ge \frac{i_j}{I_j}\ge\frac{1}{j}$。若原调和和有限，则分块和有限，却支配同一调和和，并且首块因至少含两项而存在固定严格余量；可在有限和层面保留该固定余量后取极限，产生矛盾。\par\textbf{方法点评} 积分夹逼刻画尾和极限，非减分块则给出反证结构。\par\textbf{变式与易错点} 不要把“一个发散和大于另一个发散和”当作矛盾，必须先假设有限。\end{example}

\begin{example}{竞赛拓展｜Erdős素数倒数的计数证明}{难度：$\star\star\star$}\label{deep:ex:series_source_deep:prime}\textbf{题目} 证明$\sum_{p\ {\rm prime}}\frac{1}{p}$发散，从而调和级数发散。\par\textbf{思路} 若大素数倒数尾和小，则大多数整数只含有限个小素数，然而这种整数计数太少。\par\textbf{详细解答} 反设素数倒数和收敛，选小素数$p_1,\ldots,p_K$使其余素数倒数和$<\frac{1}{2}$。$[1,N]$内至少有一个大素因子的整数个数$\le\sum_{p>p_K}\lfloor \frac{N}{p}\rfloor<\frac{N}{2}$。仅含小素数的整数可唯一写为$ab^2$，其中$a$为这些$K$个素数的无平方因子乘积，有至多$2^K$种，且$b\le\sqrt N$，故其数目$\le2^K\sqrt N$。于是$N<\frac{N}{2}+2^K\sqrt N$，当$N>2^{2K+2}$时矛盾。\par\textbf{方法点评} 核心是对两类整数分别计数，无须循环引用调和级数发散。\par\textbf{变式与易错点} 对“有某大素因子”的集合用并集上界即可，不要求互不相交。\end{example}

\begin{example}{竞赛拓展｜Erdős第二种筛除证明}{难度：$\star\star\star$}\label{deep:ex:series_source_deep:primesecond}\textbf{题目} 用无平方因子整数给出另一种素数倒数发散证明。\par\textbf{思路} 去掉被素数平方整除的数，再去掉有大素因子的数。\par\textbf{详细解答} 有$\sum_p\frac{1}{p^2}\le\frac{1}{4}+\sum_{n=3}^\infty\frac{1}{n(n-1)}=\frac{3}{4}$。若$\sum_p\frac{1}{p}$收敛，选$K$使大素数尾和$<\frac{1}{8}$。$[1,N]$中非无平方因子数至多$\frac{3N}{4}$，有大素因子的数少于$\frac{N}{8}$。因此至少$\frac{N}{8}$个整数既无平方因子又只含$K$个小素数，而这样的整数最多$2^K$个。令$N>8\cdot2^K$矛盾。\par\textbf{方法点评} 本证明直接筛除不满足条件的整数，与平方部分分解形成互补。\par\textbf{变式与易错点} 这里的计数只需上界，交叠删除不会损害下界。\end{example}

\begin{example}{竞赛拓展｜控制收敛反证与无界区间}{难度：$\star\star\star$}\label{deep:ex:series_source_deep:dominated}\textbf{题目} 通过函数积分证明调和级数发散，并解释无界区间上一致收敛为何不足以交换积分。\par\textbf{思路} 反设调和和有限可构造可积控制函数。\par\textbf{详细解答} 令$g(x)=\frac{1}{k}$在$k-1\le x<k$上，若调和和有限，则$g$可积。令$f_n(x)=n^{-1}\mathbf1_{[0,n)}(x)$。逐点$f_n\to0$，且当$x\in[k-1,k)$、$x\le n$时$k\le n$，所以$f_n\le\frac{1}{k}=g(x)$。控制收敛会给$\int f_n\to0$，但每个积分均为1，矛盾。又$\|f_n\|_\infty=\frac{1}{n}\to0$，说明在无界区间仅一致收敛不能保证积分收敛。\par\textbf{方法点评} 控制收敛反证同时揭示了无界区间上交换积分与极限所需的额外条件。\par\textbf{变式与易错点} 有限区间的一致收敛积分定理不能直接搬到无限测度区间。\end{example}

\begin{example}{竞赛风格｜均值不等式的有限和反证}{难度：$\star\star\star$}\label{deep:ex:series_source_deep:means}\textbf{题目} 用成对均值不等式证明调和级数发散。\par\textbf{思路} 假设总和有限，比较每对$(2n+1,2n+2)$与压缩后的项。\par\textbf{详细解答} 对$a,b>0$，$\frac{1}{a}+\frac{1}{b}\ge\frac{4}{a+b}$。因此$\frac{1}{2n+1}+\frac{1}{2n+2}>\frac{1}{n+1}$。若总和为$S<\infty$，按对分组后与$\sum_{n\ge0}\frac{1}{n+1}=S$比较；第一对已产生固定严格差，剩余每对差为正，所以$S>S$，矛盾。用$\frac{1}{a}+\frac{1}{b}\ge\frac{2}{\sqrt{ab}}>\frac{2}{b}$（$a<b$）也给同样的逐对严格比较。\par\textbf{方法点评} 成对比较中，严格差来自固定首对，取极限时仍被保留。\par\textbf{变式与易错点} 若不先假设$S$有限，$\infty>\infty$没有意义。\end{example}

\begin{warning}用图像、弧长或非负矩阵表达分组与面积比较时，其严谨核心仍须写为有限和不等式后取极限。埃及分数树的行内互异性质需要额外组合引理，不能仅由$\frac{1}{n}=\frac{1}{n+1}+\frac{1}{n(n+1)}$一句推出。预设Riemann假设的推理仅为条件推理，不能代替无条件证明；调和级数发散已由前面的初等证明建立。\end{warning}

\subsection{幂级数求和的微分、积分与方程法}\label{deep:series_source_deep:summing}

\begin{method}{幂级数求和的微分、积分与方程法}\textbf{适用条件} 在收敛半径内部可逐项运算，端点另证。\par\textbf{步骤与依据} 观察系数中的$n,n+1,2n-1$，分别选择导数、积分、奇偶筛选或ODE。\par\textbf{易错点} 端点值不能只由内部闭式硬代入。\end{method}

\begin{example}{考研风格｜带$n^2$的和三解}{难度：$\star\star\star$}\label{deep:ex:series_source_deep:nsquare}\textbf{题目} 求$\sum_{n=1}^\infty n^2x^n$，$|x|<1$。\par\textbf{思路} 微分几何级数、差分系数或分解$n^2=n(n-1)+n$。\par\textbf{详细解答} \textbf{微分法：}$\sum nx^n=\frac{x}{(1-x)^2}$，再作用$xD$得$\frac{x(1+x)}{(1-x)^3}$。\par\textbf{错位相减：}设$S=\sum n^2x^n$，则$(1-x)S=\sum(2n-1)x^n=\frac{2x}{(1-x)^2}-\frac{x}{1-x}$，解出同式。\par\textbf{组合导数：}$\sum n(n-1)x^n=\frac{2x^2}{(1-x)^3}$，加$\frac{x}{(1-x)^2}$。\par\textbf{方法点评} 微分法系统，错位法可以只用绝对收敛的级数代数。\par\textbf{变式与易错点} $x=\pm1$处级数均不收敛。\end{example}

\begin{example}{考研风格｜分母结构与逐项求导求和}{难度：$\star\star\star$}\label{deep:ex:series_source_deep:jdenominator}\textbf{题目} 求$\sum_{n=1}^\infty\frac{(-1)^{n-1}}{n(2n-1)}$。\par\textbf{思路} 引入$x^{2n}$后求导两次，或用积分表示。\par\textbf{详细解答} 设$S(x)=\sum\frac{(-1)^{n-1}x^{2n}}{n(2n-1)}$，$|x|<1$。则$S''=\frac{2}{1+x^2}$，由$S(0)=S'(0)=0$得$S=2x\arctan x-\ln(1+x^2)$。原级数绝对收敛，故由一致收敛可令$x\uparrow1$，得$\frac{\pi}{2}-\ln2$。另一法$\frac{1}{n(2n-1)}=\frac{2}{2n-1}-\frac{1}{n}$，两交错和分别为$\frac{\pi}{4}$和$\ln2$，线性组合也得结果。\par\textbf{方法点评} 第二法更短，第一法展示“分母变导数”的通用思路。\par\textbf{变式与易错点} 拆分后的两个级数均收敛才可分别求和。\end{example}

\subsection{一致收敛与逐项运算的反例}\label{deep:series_source_deep:uniform}

\begin{method}{一致收敛与逐项运算的反例}\textbf{适用条件} 函数列或函数项级数定义在同一集合上。\par\textbf{步骤与依据} 使用尾和上确界、M判别或一致Cauchy准则；求导需导数级数一致收敛及一点原级数收敛。\par\textbf{易错点} 逐点收敛不能自动保持连续、积分或导数。\end{method}

\begin{example}{考研风格｜一致收敛的域依赖}{难度：$\star\star$}\label{deep:ex:series_source_deep:uniform1}\textbf{题目} 讨论$f_n(x)=x^n$在$[0,1]$和$[0,a]$（$a<1$）上的一致收敛。\par\textbf{思路} 计算误差上确界，比较极限连续性。\par\textbf{详细解答} 在$[0,1]$，极限$f(x)=0$（$x<1$）、$f(1)=1$，不连续，故非一致；直接有$\sup_{x<1}x^n=1$。在$[0,a]$，$\|f_n\|_\infty=a^n\to0$，所以一致。若在$[0,1)$，极限处处0，但上确界仍为1，仍不一致。\par\textbf{方法点评} 删去一个端点不一定改善一致性，关键在是否远离该端点。\par\textbf{变式与易错点} 上确界不一定取得，不能误写最大值必存在。\end{example}

\begin{example}{竞赛风格｜导数极限不可自动交换}{难度：$\star\star\star$}\label{deep:ex:series_source_deep:derivative}\textbf{题目} 构造$f_n\to0$一致而$f_n'$不收敛到0的例子，并陈述正确的充分条件。\par\textbf{思路} 小振幅配高频率。\par\textbf{详细解答} 取$f_n(x)=\frac{\sin(nx)}{n}$，在$\R$上一致趋0，而$f_n'(0)=1$，故导数不趋0。正确条件：若$f_n'$在$[a,b]$上一致趋$g$，且$f_n(x_0)$收敛，则$f_n(x)=f_n(x_0)+\int_{x_0}^xf_n'$一致收敛到$f(x)=c+\int_{x_0}^xg$，若$g$连续则$f'=g$。\par\textbf{方法点评} 由微积分基本定理证明，比单独背交换条件更可靠。\par\textbf{变式与易错点} 仅导数逐点收敛仍不足以给上述统一积分估计。\end{example}

\subsection{Abel端点定理与Raabe临界精化}\label{deep:series_source_deep:abelproof}

\begin{method}{Abel端点定理与Raabe临界精化}\textbf{适用条件} 普通数项级数已收敛，或正项比值具有一阶渐近。\par\textbf{步骤与依据} Abel用部分和表示；Raabe与$p$级数比较，比较时选严格介于1与极限之间的指数。\par\textbf{易错点} 端点定理不能从内部闭式反推原级数一定收敛。\end{method}

\begin{theorem}{Abel端点定理的证明}若$\sum_{n=0}^\infty a_n=S$，则$\lim_{r\uparrow1}\sum_{n=0}^\infty a_nr^n=S$。\end{theorem}\begin{proof}令$S_n=\sum_{k=0}^na_k$，分部求和给$\sum a_nr^n=(1-r)\sum S_nr^n$，其中$S_n$有界保证右边绝对收敛。减去$S=(1-r)\sum Sr^n$，误差为$(1-r)\sum(S_n-S)r^n$。固定$N$，首段有限项乘$1-r$趋0；尾段$|S_n-S|<\varepsilon$，其加权和$\le\varepsilon$。先$r\uparrow1$再$\varepsilon\downarrow0$。\end{proof}

\begin{theorem}{Raabe判别的严格比较证明}$a_n>0$且$n(\frac{a_n}{a_{n+1}}-1)\to r$。若$r>1$则$\sum a_n$收敛；若$r<1$则发散；$r=1$不定。\end{theorem}\begin{proof}若$r>1$，取$1<p<r$。取$\eta=\frac{r-p}{2}>0$，由极限定义，最终$\frac{a_n}{a_{n+1}}\ge1+\frac{p+\eta}{n}$；而$(1+\frac{1}{n})^p=1+\frac{p}{n}+O(n^{-2})\le1+\frac{p+\eta}{n}$最终成立。两者严格比较，得到$a_{n+1}(n+1)^p\le a_nn^p$最终成立，所以$a_n\le Cn^{-p}$，收敛。若$r<1$，取$r<p<1$，反向得到$a_nn^p$最终递增，故$a_n\ge cn^{-p}$，发散。$r=1$的例子$\frac{1}{n}$发散、$\frac{1}{n(\ln n)^2}$收敛。\end{proof}

\begin{example}{竞赛风格｜Abel和存在不等于普通收敛}{难度：$\star\star$}\label{deep:ex:series_source_deep:endpointtrap}\textbf{题目} 解释为什么$\sum_{n\ge0}(-1)^nr^n=\frac{1}{1+r}$在$r\uparrow1$趋$\frac{1}{2}$，却不能说$\sum(-1)^n=\frac{1}{2}$。\par\textbf{思路} Abel定理只有从普通收敛到径向极限的方向。\par\textbf{详细解答} 当$r<1$，几何级数绝对收敛。径向极限$\frac{1}{2}$称为该发散级数的Abel和。原级数部分和在1和0间跳动，没有普通和。Abel端点定理的前提$\sum a_n$已收敛在此不成立，故不能逆用。\par\textbf{方法点评} 复数求和及径向极限都必须补充端点收敛的逻辑。\par\textbf{变式与易错点} 使用广义求和时必须明确标记，不可与普通和混写。\end{example}

\section{Fourier点收敛的基础证明与进阶接口}\label{deep:fourier_source_deep}

\begin{examtip}本节展开核、收敛、极大算子与函数空间的基本证明。Carleson极大估计作为高级标准结论使用；以下基础引理说明由主估计通向收敛结论的逻辑。\end{examtip}

\subsection{Riemann--Lebesgue引理与局部化}\label{deep:fourier_source_deep:riemann}

\begin{method}{Riemann--Lebesgue引理与局部化}\textbf{适用条件} $f\in L^1$，利用简单函数或阶梯函数的$L^1$逼近。\par\textbf{步骤与依据} 先对区间示性函数直接积分，再用$|\widehat f-\widehat g|\le\|f-g\|_1$推广。\par\textbf{易错点} 连续性不是必要条件；可积性和逼近范数才是关键。\end{method}

\begin{theorem}{Riemann--Lebesgue引理}若$f\in L^1([-\pi,\pi])$，则$\widehat f(n)\to0$。\end{theorem}\begin{proof}对区间$[a,b]$，$\int_a^b\e^{-int}\dd t=\frac{\e^{-ina}-\e^{-inb}}{in}\to0$，故有限阶梯函数成立。给$\varepsilon>0$取阶梯函数$g$使$\|f-g\|_1<\varepsilon$，则$2\pi|\widehat f(n)|\le\|f-g\|_1+|\int g\e^{-int}|\le\varepsilon+o(1)$，令$\varepsilon\downarrow0$。\end{proof}

\begin{theorem}{Riemann局部化原理}若两个周期可积函数在$x$的某邻域相同，则它们的Fourier部分和在$x$之差趋0。\end{theorem}\begin{proof}令$h=f-g$在$|t-x|<\delta$为0。由Dirichlet核，$S_Nh(x)$是远离0区间上$\frac{h(x-u)}{\sin(\frac{u}{2})}$与$\sin((N+\frac{1}{2})u)$的积分。该商可积，Riemann--Lebesgue引理给极限0。半整数频率可吸收固定的$\sin(\frac{u}{2}),\cos(\frac{u}{2})$因子，仍由整数频率的引理处理。\end{proof}

\begin{example}{竞赛拓展｜Dini收敛准则的完整核推导}{难度：$\star\star\star$}\label{deep:ex:fourier_source_deep:dini}\textbf{题目} 设$f$在$x$两侧极限为$f_+,f_-$，且$\int_0^\delta\frac{|f(x+t)+f(x-t)-f_+-f_-|\dd t}{t}<\infty$。证明$S_Nf(x)\to\frac{f_++f_-}{2}$。\par\textbf{思路} 核的偶性把两侧值成对，并将误差写成可积函数的振荡积分。\par\textbf{详细解答} 设$s=\frac{f_++f_-}{2}$。常函数核积分为1。记
\[g(t)=\dfrac{f(x+t)+f(x-t)-2s}{\sin(\frac{t}{2})},\]
则误差可写为
\[S_Nf(x)-s=\dfrac{1}{2\pi}\int_0^\pi
     g(t)\sin\bigl((N+\dfrac12)t\bigr)\dd t.\]
近0时，$g$由假设可积；远离0时，由$f\in L^1$可知$g$可积。Riemann--Lebesgue引理使整体趋0。\par\textbf{方法点评} 给出的充分条件是对称差的Dini条件，比单侧绝对差条件稍灵活。\par\textbf{变式与易错点} 若只知道$f$连续，商未必可积，不能直接沿用此证明。\end{example}

\begin{example}{考研风格｜跳跃点、端点与奇偶延拓}{难度：$\star\star\star$}\label{deep:ex:fourier_source_deep:jump}\textbf{题目} 对$f(x)=x$在$(-\pi,\pi)$的周期延拓，写Fourier级数并说明端点值。\par\textbf{思路} 奇函数只有正弦项，跳跃点收敛到左右平均。\par\textbf{详细解答} $b_n=(\frac{2}{\pi})\int_0^\pi x\sin nx\dd x=\frac{2(-1)^{n+1}}{n}$。所以$x=2\sum_{n\ge1}\frac{(-1)^{n+1}\sin nx}{n}$在$-\pi<x<\pi$成立。在$x=\pi$和$-\pi$，周期左右极限为$\pi,-\pi$，平均为0，而每个正弦项也为0。\par\textbf{方法点评} 区间端点按周期延拓处理，不能只看原区间单侧。\par\textbf{变式与易错点} 给$f(\pi)$指定成$\pi$不改变系数，也不改变级数在该点收敛到0。\end{example}

\subsection{Fejér核、逼近与Fourier唯一性}\label{deep:fourier_source_deep:fejerproof}

\begin{method}{Fejér核、逼近与Fourier唯一性}\textbf{适用条件} 周期连续函数，或$L^p$中的平移连续性。\par\textbf{步骤与依据} 核非负、质量1、远端质量趋0三个性质构成近似恒等算子。\par\textbf{易错点} Cesàro平均与普通部分和是不同序列。\end{method}

\begin{theorem}{Fejér核的三条性质}$F_N(t)=\frac1{N+1}\left|\sum_{k=0}^N\e^{ikt}\right|^2$非负，$(2\pi)^{-1}\int F_N=1$，且对每个$0<\delta\le\pi$，$\sup_{\delta\le|t|\le\pi}F_N(t)\le[(N+1)\sin^2(\frac{\delta}{2})]^{-1}$。\end{theorem}\begin{proof}非负由平方给出。展开平方$\frac1{N+1}\sum_{j,k=0}^N\e^{i(j-k)t}$，只有$j=k$项积分非零，恰有$N+1$项。几何和公式给$F_N=\frac{(N+1)^{-1}\sin^2(\frac{(N+1)t}{2})}{\sin^2(\frac{t}{2})}$，分母远离0即得尾界。\end{proof}

\begin{example}{竞赛拓展｜Fourier系数唯一性}{难度：$\star\star\star$}\label{deep:ex:fourier_source_deep:uniqueness}\textbf{题目} 若周期连续函数所有Fourier系数均为0，证明函数恒为0。\par\textbf{思路} 所有Fejér均值为0，而它们一致趋原函数。\par\textbf{详细解答} 由系数表达，$\sigma_Nf=\sum_{|k|\le N}(1-\frac{|k|}{N+1})\widehat f(k)\e^{ikx}=0$。另一方面核三性质和一致连续性给$\|\sigma_Nf-f\|_\infty\to0$，故$f=0$。对$L^1$函数，用平移的$L^1$连续性证明$\|\sigma_Nf-f\|_1\to0$，可得$f=0$几乎处处。\par\textbf{方法点评} 唯一性不要求普通Fourier部分和已逐点收敛。\par\textbf{变式与易错点} $L^1$版本只能得到几乎处处相同，不能识别零测集上的任意改值。\end{example}

\subsection{弱型估计、分布函数与强型推导}\label{deep:fourier_source_deep:maxproof}

\begin{method}{弱型估计、分布函数与强型推导}\textbf{适用条件} 实轴极大算子先定义在$L^1_{\rm loc}$上。\par\textbf{步骤与依据} 用有限区间覆盖选不交子族；按每个阈值重新拆分函数；最后积分分布函数。\par\textbf{易错点} 从每个单独平均的有界性不能直接得到上确界算子的同型界。\end{method}

\begin{theorem}{从弱$(1,1)$到强$L^p$的直接推导}若次线性正极大算子$M$满足$|\{Mf>\lambda\}|\le \frac{C\|f\|_1}{\lambda}$及$\|Mf\|_\infty\le\|f\|_\infty$，则对$p>1$有$\|Mf\|_p^p\le \frac{C\,2^p p\|f\|_p^p}{p-1}$。\end{theorem}\begin{proof}按阈值分解$f=g+h$，其中$g=f\mathbf1_{\{|f|\le\frac{\lambda}{2}\}}$。有$Mg\le\frac{\lambda}{2}$，故$\{Mf>\lambda\}\subset\{Mh>\frac{\lambda}{2}\}$，测度$\le2C\lambda^{-1}\int_{|f|>\frac{\lambda}{2}}|f|$。乘$p\lambda^{p-1}$积分并交换非负积分，得到$2Cp\int|f(x)|\int_0^{2|f(x)|}\lambda^{p-2}\dd\lambda\dd x=\frac{C2^pp\|f\|_p^p}{p-1}$。先作截断避免无穷，再用单调收敛。\end{proof}

\begin{example}{竞赛拓展｜Lebesgue微分定理的误差集证明}{难度：$\star\star\star$}\label{deep:ex:fourier_source_deep:differentiation}\textbf{题目} 说明如何由极大函数弱型与连续函数稠密性，证明$f\in L^1$几乎每点为Lebesgue点。\par\textbf{思路} 以连续近似$g$分拆局部平均振幅。\par\textbf{详细解答} 记$D f(x)=\limsup_{r\downarrow0}(2r)^{-1}\int_{x-r}^{x+r}|f(t)-f(x)|\dd t$。连续$g$满足$Dg=0$，故$Df\le M(f-g)+|f-g|$。于是$\lvert\{Df>2\varepsilon\}\rvert\le\dfrac{(C+1)\lVert f-g\rVert_1}{\varepsilon}$，第二项用Chebyshev。取连续紧支撑$g$使$L^1$误差任意小，则该集合测度为0。对可数$\varepsilon=\frac{1}{m}$合并，得$Df=0$几乎处处。\par\textbf{方法点评} 误差拆成极大平均与点值两项，补齐了“用稠密性推广”中的关键估计。\par\textbf{变式与易错点} 局部可积情形先在有界区间截断，再用可数区间覆盖全实轴。\end{example}

\subsection{奇异积分的抵消、频率乘子与主值}\label{deep:fourier_source_deep:hilbertproof}

\begin{method}{奇异积分的抵消、频率乘子与主值}\textbf{适用条件} 先对光滑紧支撑函数定义，再借范数控制延拓。\par\textbf{步骤与依据} 局部减去常数；远端直接积分；用Fourier乘子看$L^2$结构。\par\textbf{易错点} 主值存在与两个单侧反常积分收敛不是同一件事。\end{method}

\begin{example}{竞赛拓展｜示性函数的Hilbert变换}{难度：$\star\star\star$}\label{deep:ex:fourier_source_deep:hilbertindicator}\textbf{题目} 计算$H\mathbf1_{[0,1]}(x)$，采用$Hf(x)=\pi^{-1}\operatorname{p.v.}\int \frac{f(t)}{x-t}\dd t$。\par\textbf{思路} 直接积分对数，区间内用对称截断抵消。\par\textbf{详细解答} 当$x\notin[0,1]$，积分$=\pi^{-1}[\ln|x|-\ln|x-1|]$。当$0<x<1$，在$t=x$左右截去同一$\varepsilon$，两项中的$\ln\varepsilon$抵消，得到同一表达式。因此除端点外$Hf(x)=\pi^{-1}\ln|\frac{x}{x-1}|$。端点对数发散；远端$Hf(x)\sim\frac{1}{\pi x}$。\par\textbf{方法点评} 给出主值、端点奇性和弱型尾部的具体模型。\par\textbf{变式与易错点} 若采用核$\frac{1}{t-x}$，整体符号相反。\end{example}

\begin{example}{竞赛拓展｜Hilbert变换的$L^2$等距}{难度：$\star\star\star$}\label{deep:ex:fourier_source_deep:l2hilbert}\textbf{题目} 在本册实轴Fourier约定下解释$H$的乘子及$L^2$性质。\par\textbf{思路} 对称截断核的Fourier变换趋$-i\operatorname{sgn}\xi$。\par\textbf{详细解答} 核$\pi^{-1}\operatorname{p.v.}(\frac{1}{x})$的乘子为$-i\operatorname{sgn}\xi$：奇性消去余弦，正弦积分
\[
\int_0^\infty\dfrac{\sin(2\pi\xi x)}x\dd x
=\dfrac\pi2\operatorname{sgn}\xi\qquad(\xi\ne0)
\]
给出该乘子。由本书归一化下的Plancherel恒等式，
\[
\|Hf\|_2^2
=\int_\R|{-i\operatorname{sgn}\xi}|^2|\widehat f(\xi)|^2\dd\xi
=\|f\|_2^2.
\]又乘子平方为$-1$几乎处处，所以$H^2f=-f$。\par\textbf{方法点评} 这是$L^2$空间中的定义与延拓；一般逐点主值仍需极大截断理论。\par\textbf{变式与易错点} 零频率的一个点不影响实轴$L^2$等式；周期情形常数项会被消去。\end{example}

\subsection{极大控制如何真正推出点收敛}\label{deep:fourier_source_deep:density}

\begin{method}{极大控制如何真正推出点收敛}\textbf{适用条件} $T_n$在稠密子类上逐点收敛，$T^*=\sup_n|T_n|$有弱型控制。\par\textbf{步骤与依据} 定义尾部振幅，近似函数的尾部振幅为0；以极大估计控制余项。\par\textbf{易错点} 范数收敛、依测度收敛与几乎处处收敛应分别证明。\end{method}

\begin{theorem}{稠密类推广引理}设$T_n$为线性算子，$1\le p<\infty$，$|\{T^*h>\lambda\}|\le C\lambda^{-p}\|h\|_p^p$，且$T_n g$在稠密类$\mathcal D\subset L^p$上几乎处处收敛，则$T_nf$对任意$f\in L^p$几乎处处为Cauchy列。\end{theorem}\begin{proof}令$\Omega(f)=\lim_{N\to\infty}\sup_{m,n\ge N}|T_mf-T_nf|$。线性给$\Omega(f)\le\Omega(g)+2T^*(f-g)=2T^*(f-g)$几乎处处。于是$|\{\Omega(f)>\varepsilon\}|\le C(\frac{2}{\varepsilon})^p\|f-g\|_p^p$。稠密性使右端任意小，故$\Omega(f)=0$几乎处处。若另知$T_nf\to f$于$L^p$或依测度，可由子列唯一性识别极限。\end{proof}

\begin{example}{竞赛拓展｜Carleson主估计的逻辑位置}{难度：$\star\star\star$}\label{deep:ex:fourier_source_deep:carlesongap}\textbf{题目} 从$\|S^*f\|_p\le C_p\|f\|_p$推导$L^p$ Fourier级数几乎处处收敛，$p>1$，并说明尚需证明的核心。\par\textbf{思路} 三角多项式作稠密类，套推广引理。\par\textbf{详细解答} 三角多项式$P$的$S_NP$最终等于$P$。强型估计经Chebyshev给推广引理所需弱型，故$S_Nf$几乎处处收敛。又$S_N$在$L^p$一致有界，利用多项式逼近得$\|S_Nf-f\|_p\to0$，于是点态极限为$f$。真正困难是最开始的$\|S^*f\|_p$估计，固定频率Hilbert有界性不能替代它；建立该估计需要局部范数、多尺度分解、频率配对和例外集控制。\par\textbf{方法点评} 完整写出“主估计之后”的推论，准确区分基础引理与主估计。\par\textbf{变式与易错点} 本例证明的是主估计的推论，主估计本身仍作为高级标准结论使用。\end{example}

\subsection{重排、Lorentz空间与端点区别}\label{deep:fourier_source_deep:rearrangement}

\begin{method}{重排、Lorentz空间与端点区别}\textbf{适用条件} 用分布函数组织函数值大小，测度空间先取有限或$\sigma$有限。\par\textbf{步骤与依据} 定义递减重排$f^*$，用层蛋糕公式连接范数；弱$L^p$对应分布尾界。\par\textbf{易错点} $L^{p,\infty}$不等于$L^p$，而且端点嵌入可能严格。\end{method}

\begin{definition}{递减重排与Lorentz量}令$d_f(\lambda)=|\{|f|>\lambda\}|$，$f^*(t)=\inf\{\lambda\ge0:d_f(\lambda)\le t\}$。对$0<p<\infty,0<q<\infty$，定义$\|f\|_{p,q}=(\int_0^\infty\frac{[t^{\frac{1}{p}}f^*(t)]^q\dd t}{t})^{\frac{1}{q}}$，$q=\infty$时取上确界。本册采用这个常见归一化；等价范数可有固定因子。\end{definition}

\begin{example}{竞赛拓展｜弱$L^p$与强$L^p$的严格差别}{难度：$\star\star\star$}\label{deep:ex:fourier_source_deep:weaklp}\textbf{题目} 在$(0,1)$上考察$f(x)=x^{-\frac{1}{p}}$，$p>0$。\par\textbf{思路} 直接算分布函数和积分。\par\textbf{详细解答} $\lambda\ge1$时$d_f(\lambda)=\lambda^{-p}$，$0<\lambda<1$时为1。因此$\sup_{\lambda>0}\lambda d_f(\lambda)^{\frac{1}{p}}=1$，$f\in L^{p,\infty}$。但$\int_0^1|f|^p\dd x=\int_0^1\frac{\dd x}{x}=\infty$，不在$L^p$。其重排$f^*(t)=t^{-\frac{1}{p}}$（$0<t<1$），$t\ge1$时为0；有限$q$的Lorentz量包含$\int_0^1\frac{\dd t}{t}$，均发散。\par\textbf{方法点评} 端点估计需要区分强型与弱型空间，仅用普通$L^p$记号不足以表达。\par\textbf{变式与易错点} 有限测度上$L^r\subset L^p$需要$r>p$，方向不可倒置。\end{example}

\begin{warning}接近$L^1$的端点理论涉及重排空间、Q与QA空间以及Antonov端点定理。前节给出定义与标准结论，本节说明它们所依赖的分布函数、插值与稠密推广。空间之间的精细包含关系、最优常数与时频组合主估计属于进一步专题，不能由本节基础引理直接推出。\end{warning}

\subsection{正交投影、Bessel不等式与Parseval}\label{deep:fourier_source_deep:parseval}

\begin{method}{正交投影、Bessel不等式与Parseval}\textbf{适用条件} 圆周$L^2$采用归一化内积$(2\pi)^{-1}\int f\overline g$。\par\textbf{步骤与依据} 先有限投影得到平方范数恒等式，再用三角多项式稠密性。\par\textbf{易错点} 系数平方和有限是范数结论，本身不给每点收敛。\end{method}

\begin{theorem}{Bessel不等式与Parseval恒等式}对$f\in L^2$，$\sum_{k\in\mathbb Z}|\widehat f(k)|^2=(2\pi)^{-1}\int_{-\pi}^{\pi}|f|^2$，且$S_Nf\to f$于$L^2$。\end{theorem}\begin{proof}指数函数是正交单位系，有限展开给$\|f-S_Nf\|_2^2=\|f\|_2^2-\sum_{|k|\le N}|\widehat f(k)|^2\ge0$，先得Bessel不等式。三角多项式在$L^2$稠密，若$P$次数$\le N$，正交投影的最佳逼近性给$\|f-S_Nf\|_2\le\|f-P\|_2$。因此误差趋0，有限恒等式取极限得Parseval。稠密性可由Fejér逼近连续函数和连续函数在$L^2$中稠密得到。\end{proof}

\begin{example}{竞赛风格｜平方积分求四次倒数和}{难度：$\star\star\star$}\label{deep:ex:fourier_source_deep:baselparseval}\textbf{题目} 由$f(x)=x^2$的Fourier展开求$\sum_{n\ge1}n^{-4}$。\par\textbf{思路} 先求系数，再用Parseval。\par\textbf{详细解答} $a_0=\frac{2\pi^2}{3},a_n=\frac{4(-1)^n}{n^2},b_n=0$。实形式Parseval为$\pi^{-1}\int_{-\pi}^{\pi}f^2=\frac{a_0^2}{2}+\sum(a_n^2+b_n^2)$。代入$\frac{2\pi^4}{5}=\frac{2\pi^4}{9}+16\sum n^{-4}$，所以和为$\frac{\pi^4}{90}$。若用复系数归一化，正负频率各贡献一次，得到同值。\par\textbf{方法点评} Fourier求和法可系统推广到偶次倒数和。\par\textbf{变式与易错点} 常数项为$\frac{a_0}{2}$，Parseval中贡献$\frac{a_0^2}{2}$，不能漏因子。\end{example}

\subsection{连续函数也可能Fourier发散}\label{deep:fourier_source_deep:continuousdivergence}

\begin{method}{连续函数也可能Fourier发散}\textbf{适用条件} 考虑$C(\mathbb T)$上的点值部分和泛函。\par\textbf{步骤与依据} 证明Dirichlet核$L^1$范数无界，再用一致有界原理。\par\textbf{易错点} Fejér均值一致收敛不意味着普通部分和对所有连续函数收敛。\end{method}

\begin{theorem}{存在连续函数在一点的Fourier部分和无界}存在$f\in C(\mathbb T)$使$\sup_N|S_Nf(0)|=\infty$。\end{theorem}\begin{proof}令$L_N(f)=S_Nf(0)$，其算子范数为$(2\pi)^{-1}\|D_N\|_1$；上界由积分估计，下界可用连续函数逼近$\operatorname{sgn}D_N$达到任意接近。又在$t\in[\frac{1}{N},1]$上$|\sin(\frac{t}{2})|\le \frac{t}{2}$。把该区间分成$\sin((N+\frac{1}{2})t)$的半周期，每个半周期取正弦绝对值至少$\frac{1}{2}$的中间部分，其积分下界为常数乘$\frac{1}{j}$，相加得$\|D_N\|_1\ge c\log N$。因此$\sup_N\|L_N\|=\infty$。\par 若每个连续$f$的$\sup_N|L_Nf|$都有限，由一致有界原理应有$\sup_N\|L_N\|<\infty$，矛盾。该原理的Baire证明为：闭集$E_m=\{f:\sup_N|L_Nf|\le m\}$覆盖Banach空间，故某个$E_m$有内点；相减得到零点某球内全部$L_N$一致有界，缩放即得范数一致有界。\end{proof}

\begin{example}{竞赛拓展｜三种收敛互不替代}{难度：$\star\star\star$}\label{deep:ex:fourier_source_deep:threelimits}\textbf{题目} 把普通逐点收敛、几乎处处收敛与Fejér一致收敛在连续函数上的关系说清。\par\textbf{思路} 结合上一存在性定理、Carleson--Hunt和Fejér。\par\textbf{详细解答} 每个连续周期函数属于$L^2$，所以由Carleson--Hunt其普通Fourier级数几乎处处收敛。上一结论给出某个连续函数在指定点普通部分和无界，故不能提升到每点。另一方面每个连续函数的Fejér均值一致收敛。这三句相容：平均序列更平滑，且几乎处处允许例外零测集。\par\textbf{方法点评} 消除把连续性、正交性和点收敛混为一谈的常见误读。\par\textbf{变式与易错点} 例外集可非空；一致收敛必逐点收敛，但说的是同一个序列时才可推。\end{example}

\section{级数判别与求和的综合构型}\label{ic:delta_series}

\begin{examtip}本节以几何级数、裂项和逐项求导积分为基础，进一步讨论判别细化、参数族与综合算例。\end{examtip}

\subsection{对数判别与统一裕量}\label{ic:delta_series:logtest}\begin{method}{对数判别与统一裕量}\textbf{识别信号} 正项可以写为指数形式，或根值、比值都处于临界值1。\par \textbf{操作步骤} 计算$\lambda_n=-\frac{\ln a_n}{\ln n}$。若最终$\lambda_n\ge1+\delta$且$\delta>0$固定，收敛；若最终$\lambda_n\le1$，发散。极限$\lambda>1$或$<1$时分别推出上述比较。\par \textbf{疑难点与条件} 最终$\lambda_n>1$但与1的差趋0，不保证收敛。\end{method}

\begin{theorem}{对数判别的比较证明}若$\lambda_n\ge1+\delta$，则$a_n\le n^{-1-\delta}$；若$\lambda_n\le1$，则$a_n\ge\frac{1}{n}$。因此有上述敛散结论。\end{theorem}\begin{proof}由于$n>1$时$\ln n>0$，$\lambda_n\ge1+\delta$等价于$\ln a_n\le-(1+\delta)\ln n$。指数函数保序给第一界。第二界同理。再使用$p$级数与调和级数比较即可。\end{proof}

\begin{example}{竞赛风格｜指数写法隐藏的幂级数}{难度：$\star\star\star$}\label{ic:ex:delta_series:logparameter}\textbf{题目} 判定$\sum_{n\ge2}b^{-\ln n}$，$b>1$。\par \textbf{思路} 将底数幂换成自然指数。\par \textbf{详细解答} $b^{-\ln n}=\e^{-(\ln b)(\ln n)}=n^{-\ln b}$。对数判别量恒为$\ln b$，所以$b>\e$时收敛，$1<b<\e$时发散；$b=\e$时恰为调和尾和，也发散。\par \textbf{验算与条件复核} 指数恒等式严格成立，无渐近误差；阈值由$p=\ln b$与1比较确定。\par \textbf{疑难点与方法比较} 临界值需回到原级数，不可把判别失效写成敛散未必能知。\end{example}

\subsection{Raabe非极限版与临界反例}\label{ic:delta_series:raabemargin}\begin{method}{Raabe非极限版与临界反例}\textbf{识别信号} $r_n=n(\frac{a_n}{a_{n+1}}-1)$不一定有极限。\par \textbf{操作步骤} 若最终$r_n\ge1+\delta$，取$1<p<1+\delta$与$n^{-p}$作比值比较；若最终$r_n\le1$，则$na_n$最终不减，支配调和级数。\par \textbf{疑难点与条件} “$r_n>1$”与“距1有固定正裕量”不同。\end{method}

\begin{theorem}{Raabe单侧形式}最终$r_n\ge1+\delta$时收敛；最终$r_n\le1$时发散。\end{theorem}\begin{proof}第一种情形选$1<p<1+\delta$，因$(1+\frac{1}{n})^p=1+\frac{p}{n}+O(n^{-2})$，最终$\frac{a_n}{a_{n+1}}\ge(1+\frac{1}{n})^p$，故$n^pa_n$最终不增，$a_n\le Cn^{-p}$。第二种有$\frac{a_n}{a_{n+1}}\le\frac{n+1}{n}$，即$na_n\le(n+1)a_{n+1}$，所以$a_n\ge \frac{c}{n}$。\end{proof}

\begin{example}{疑难辨析｜每项判别量大于1仍可发散}{难度：$\star\star\star$}\label{ic:ex:delta_series:raabecounter}\textbf{题目} 说明$a_n=\frac{1}{n\ln n}$满足$r_n>1$但其级数发散，$n\ge2$。\par \textbf{思路} 严格不等式不等于统一裕量。\par \textbf{详细解答} $\frac{a_n}{a_{n+1}}=\frac{(1+\frac{1}{n})\ln(n+1)}{\ln n}>1+\frac{1}{n}$，故$r_n>1$。但$\int_2^\infty\frac{\dd x}{x\ln x}=\infty$，由积分判别发散。展开得$r_n=1+\frac{1}{\ln n}+O(\frac{1}{n\ln n})\to1$，恰说明缺少固定裕量。\par \textbf{验算与条件复核} 代入比值及变量$t=\ln x$的积分变换分别核验两结论。\par \textbf{疑难点与方法比较} 不能将充分条件中的固定$\delta$删除。\end{example}

\begin{example}{竞赛风格｜阶乘商的阈值与精确和}{难度：$\star\star\star$}\label{ic:ex:delta_series:factorialparameter}\textbf{题目} 判定并在收敛时求$\sum_{n\ge1}\frac{n!}{(a+1)(a+2)\cdots(a+n)}$，$a>0$。\par \textbf{思路} 比值给Raabe量，Beta积分给总和。\par \textbf{详细解答} 记项为$u_n$，则$\frac{u_n}{u_{n+1}}=\frac{n+a+1}{n+1}$，Raabe量$\frac{an}{n+1}\to a$，故$a>1$收敛、$0<a<1$发散。$a=1$时$u_n=\frac{1}{n+1}$也发散。若$a>1$，$u_n=a\int_0^1t^n(1-t)^{a-1}\dd t$，由非负性换序，总和$a\int_0^1t(1-t)^{a-2}\dd t=\frac{1}{a-1}$。\par \textbf{验算与条件复核} Beta恒等式$aB(n+1,a)=\frac{\Gamma(n+1)\Gamma(a+1)}{\Gamma(n+a+1)}$验证项表示；末积分直接换$u=1-t$为$a[\frac{1}{a-1}-\frac{1}{a}]$。\par \textbf{疑难点与方法比较} Gamma表达、判敛和换序的条件分别交代，不能先对发散参数求和。\end{example}

\begin{example}{竞赛风格｜双阶乘幂的临界参数}{难度：$\star\star\star$}\label{ic:ex:delta_series:doublefactorial}\textbf{题目} 判定$\sum_{n\ge1}[\frac{(2n-1)!!}{(2n)!!}]^s$，$s>0$。\par \textbf{思路} Raabe先给非临界值，Wallis补$s=2$。\par \textbf{详细解答} 令$b_n$为括号内比值，$\frac{b_n}{b_{n+1}}=\frac{2n+2}{2n+1}$，所以$n[(\frac{b_n}{b_{n+1}})^s-1]\to \frac{s}{2}$。$s>2$收敛，$s<2$发散。$s=2$时，由Wallis得$b_n\sim\frac{1}{\sqrt{\pi n}}$，因此$b_n^2\sim\frac{1}{\pi n}$，发散。\par \textbf{验算与条件复核} 对$(1+\frac{1}{2n+1})^s$作一阶展开得$\frac{s}{2}$；临界项用正项极限比较。\par \textbf{疑难点与方法比较} Raabe量等于1时必须继续判断，不能简单归入收敛或发散。\end{example}

\subsection{用部分和作分母的抽象正项级数}\label{ic:delta_series:partialsums}\begin{method}{用部分和作分母的抽象正项级数}\textbf{识别信号} $u_n>0,S_n=\sum_{k=1}^nu_k$。\par \textbf{操作步骤} 若$S_n\to\infty$，$\sum \frac{u_n}{S_n^\alpha}$在$\alpha>1$时收敛，$\alpha\le1$时发散；若$S_n$有有限极限，则任意实$\alpha$均收敛。\par \textbf{疑难点与条件} 处理$\alpha=1$必须容许$u_n$有很大跳跃，不能假定$\frac{u_n}{S_n}\to0$。\end{method}

\begin{theorem}{部分和分母定理的完整证明}上述结论不要求$u_n$单调。\end{theorem}\begin{proof}当$\alpha>1$，对$n\ge2$，$\frac{u_n}{S_n^\alpha}\le\int_{S_{n-1}}^{S_n}t^{-\alpha}\dd t$，总和受$\int_{S_1}^\infty t^{-\alpha}\dd t$控制。若$\sum \frac{u_n}{S_n}$收敛，记$v_n=\frac{u_n}{S_n}$，则最终$v_n\le\frac{1}{2}$，且$\ln S_n-\ln S_{n-1}=-\ln(1-v_n)\le2v_n$，累加会使$\ln S_n$有界，与$S_n\to\infty$矛盾。$\alpha<1$时最终$S_n\ge1$，各项不小于$\frac{u_n}{S_n}$。若$S_n\to S<\infty$，因$S_n$在$[u_1,S]$内，$S_n^{-\alpha}$有界，与$\sum u_n$比较即可。\end{proof}

\begin{example}{竞赛风格｜非线性递推与级数的等价}{难度：$\star\star\star$}\label{ic:ex:delta_series:recursivesum}\textbf{题目} $b_1=1,a_n>0,b_{n+1}=b_n+\frac{a_n}{b_n}$，证明$\sum a_n$收敛当且仅当$b_n$收敛。\par \textbf{思路} 把$a_n$还原成$b_n$乘增量。\par \textbf{详细解答} $b_n\ge1$且递增。若$\sum a_n<\infty$，则$b_{n+1}-b_n=\frac{a_n}{b_n}\le a_n$，累加知$b_n$有界。反之若$b_n\to B$，则$a_n=b_n(b_{n+1}-b_n)\le B(b_{n+1}-b_n)$，部分和$\le B(B-1)$，故收敛。\par \textbf{验算与条件复核} 只使用正性、单调有界和有限望远镜求和，无未授权极限交换。\par \textbf{疑难点与方法比较} 递推中的分母$b_n$非零由$b_n\ge1$先保证。\end{example}

\begin{example}{竞赛风格｜取整函数控制的分组交错}{难度：$\star\star\star$}\label{ic:ex:delta_series:floorsign}\textbf{题目} 判定$\sum_{n\ge1}\frac{(-1)^{\lfloor\sqrt n\rfloor}}{n}$。\par \textbf{思路} 按平方区间分组，不直接声称原项交错。\par \textbf{详细解答} 令$B_k=\sum_{n=k^2}^{(k+1)^2-1}\frac{1}{n}$。由$\frac{2k+1}{(k+1)^2}\le B_k\le\frac{2k+1}{k^2}$，得到$B_k=\frac{2}{k}+O(k^{-2})$。所以$\sum(-1)^kB_k$是收敛交错调和级数加绝对收敛余项。组内任意部分的绝对值不超过$B_k\to0$，故原部分和也收敛。绝对值级数为调和级数，因而条件收敛。\par \textbf{验算与条件复核} 两界与$\frac{2}{k}$的差均为$O(k^{-2})$，严格控制组内误差。\par \textbf{疑难点与方法比较} 组长无界时不能只凭分组后收敛推出原级数收敛，必须补组内振幅趋零。\end{example}

\begin{example}{竞赛风格｜连乘分母的裂项}{难度：$\star\star\star$}\label{ic:ex:delta_series:productsum}\textbf{题目} $a>1$，求$\sum_{n\ge1}\frac{a^n}{\prod_{k=1}^n(1+a^k)}$。\par \textbf{思路} 设连乘倒数$b_n$。\par \textbf{详细解答} $b_0=1,b_n=\prod_{k=1}^n(1+a^k)^{-1}$，则$b_{n-1}-b_n=a^nb_n$。有限部分和$1-b_N$。又$0<b_N\le a^{-\frac{N(N+1)}{2}}\to0$，故总和为1。\par \textbf{验算与条件复核} 相邻乘积仅差因子$1+a^n$，所以裂项恒等式严格成立。\par \textbf{疑难点与方法比较} 先写有限和，避免不必要的无穷乘积交换。\end{example}

\begin{example}{考研风格｜复合指数的系数求和}{难度：$\star\star\star$}\label{ic:ex:delta_series:weightedexp}\textbf{题目} 求$\sum_{n\ge0}\frac{(2n+1)x^{2n}}{n!}$。\par \textbf{思路} 它是$x\e^{x^2}$的导数。\par \textbf{详细解答} 由$\sum \frac{x^{2n+1}}{n!}=x\e^{x^2}$，在任意紧区间逐项求导，得和$(1+2x^2)\e^{x^2}$。收敛半径为无穷，所以全实轴有效。\par \textbf{验算与条件复核} 系数乘$2n+1$正对应幂$x^{2n+1}$的导数，没有遗漏$n=0$项。\par \textbf{疑难点与方法比较} 逐项求导依赖幂级数的局部一致性质，而不是“各项相互独立”。\end{example}

\section{图解：部分和与平均部分和}
\begin{figure}[htbp]\centering
\begin{tikzpicture}\begin{axis}[width=.84\textwidth,height=6cm,axis lines=middle,xmin=-3.1416,xmax=3.1416,ymin=-3.6,ymax=3.6,xlabel=$x$,ylabel=$y$,samples=180,legend pos=north west]
\addplot[Ink,dashed,domain=-3.1415:3.1415]{x};\addlegendentry{周期锯齿的一个周期}
\addplot[Rust,domain=-3.1415:3.1415]{2*(sin(deg(x))-sin(deg(2*x))/2+sin(deg(3*x))/3-sin(deg(4*x))/4+sin(deg(5*x))/5-sin(deg(6*x))/6)};\addlegendentry{$S_6$}
\addplot[Teal,thick,domain=-3.1415:3.1415]{2*(6/7*sin(deg(x))-5/7*sin(deg(2*x))/2+4/7*sin(deg(3*x))/3-3/7*sin(deg(4*x))/4+2/7*sin(deg(5*x))/5-1/7*sin(deg(6*x))/6)};\addlegendentry{$\sigma_6$}
\end{axis}\end{tikzpicture}
\caption{$f(x)=x$在$(-\pi,\pi)$的周期延拓，普通部分和与Fejér平均的对照。端点跳跃处两侧平均为0；普通部分和在跳跃附近振荡，平均权重$1-\frac{|k|}{7}$抑制高频。图展示有限阶效果，不能代替点收敛定理。}\label{fig:review:fourier}
\end{figure}

\section{系数运算、隐式项与相邻消去}\label{cse:section}
\subsection{收敛半径的相消与稀疏系数}\label{cse:radius}
\begin{method}{将半径问题还原为系数或函数之差}
\textbf{步骤} 相加半径至少为两半径的最小值；若两半径不同，反设相加后半径更大，再相减产生矛盾。系数平方使用根值上极限；抽取偶系数使用$\frac{F(x)+F(-x)}{2}$。
\par\textbf{疑难点} 同半径相加可能完全抵消；子级数收敛不能反推整级数收敛。
\end{method}
\begin{theorem}{系数运算的半径规则}\label{cse:thm:radius}
$A(x)=\sum a_nx^n$、$B(x)=\sum b_nx^n$的半径分别为$R_A,R_B$。其和的半径$R_C\ge\min(R_A,R_B)$；若$R_A\ne R_B$，等号成立。
当$a_n$为实数且$0<R_A<\infty$时，$\sum a_n^2x^n$的半径为$R_A^2$。
\end{theorem}
\begin{proof}
在公共半径内部两级数绝对收敛，故其和也绝对收敛。若$R_A<R_B$而$R_C>R_A$，则$A=C-B$在半径$\min(R_C,R_B)>R_A$内绝对收敛，矛盾。
再由
$\limsup|a_n^2|^{\frac{1}{n}}=(\limsup|a_n|^{\frac{1}{n}})^2=R_A^{-2}$，
使用Cauchy--Hadamard公式即得平方系数的半径。这里“平方系数”不同于和函数$A(x)^2$。
\end{proof}
\begin{example}{考研风格｜相消后半径可以任意扩大}{难度：$\star\star$}\label{cse:ex:cancellation}
\textbf{题目} 已知$\sum a_nx^n$半径为1。$\sum(a_n+1)x^n$的半径一定等于1吗？给出两种不同结果。
\par\textbf{思路} 常数系数级数也具有半径1，但相加时可消去主系数。
\par\textbf{详细解答} 公共半径结论只保证新半径不小于1。若$a_n=1$，新系数为2，半径仍为1。
给定$s>1$，取$a_n=-1+s^{-n}$，则$|a_n|^{\frac{1}{n}}\to1$，原半径为1；相加后系数$s^{-n}$，新半径为$s$。
若$a_n=-1$，相加后恒零，半径为$\infty$。
\par\textbf{方法点评} 通过构造系数中大小不同的两层指数尺度，可精确控制相消后的半径。
\par\textbf{疑难点} 同半径相加只有下界，不能机械使用“取最小值等号”。
\end{example}
\begin{example}{竞赛风格｜偶系数的加权提取}{难度：$\star\star$}\label{cse:ex:even}
\textbf{题目} $\ln(2+x)=\sum_{n\ge0}a_nx^n$在0附近展开，求$\sum_{n\ge1}n a_{2n}$。
\par\textbf{思路} 对偶部分求导，再乘变量恢复权重。
\par\textbf{详细解答} 展开半径为2，故$x=\pm1$严格位于内部，逐项求导合法。
设$F(x)=\ln(2+x)$，则
\[
E(x)=\dfrac{F(x)+F(-x)}2=\sum_{n\ge0}a_{2n}x^{2n},
\qquad
\sum_{n\ge1}na_{2n}x^{2n}=\dfrac{xE'(x)}2
=\dfrac{x[F'(x)-F'(-x)]}{4}.
\]
令$x=1$得$\frac14(\frac{1}{3}-1)=-\frac{1}{6}$。
第二法由$a_k=\frac{(-1)^{k-1}}{k2^k}$（$k\ge1$）直接得到$na_{2n}=-\frac{1}{2\cdot4^n}$，几何求和同得$-\frac{1}{6}$。
\par\textbf{方法点评} 函数筛选法不必先写出全部系数，适合奇偶索引与导数权重同时出现。
\par\textbf{疑难点} 链式法则使$F(-x)$求导带负号；$a_0$不参与加权和。
\end{example}

\subsection{不依赖单调性的相邻消去}\label{cse:pair}
\begin{method}{有限部分和先配对}
\textbf{条件} 相邻项以同一结构出现，或递推允许还原为差分。
\par\textbf{步骤} 写出前几项，保留首尾项，证明中间项逐一抵消；普通收敛和绝对收敛分别判定。
\par\textbf{疑难点} 一般项近似于交错调和项不足以判敛；配对恒等式可以替代缺失的单调性。
\end{method}
\begin{example}{考研风格｜相邻倒数和的条件收敛}{难度：$\star\star\star$}\label{cse:ex:paired}
\textbf{题目} $u_n\ne0$且$\frac{n}{u_n}\to1$，证明
$\sum_{n\ge1}(-1)^{n-1}(\frac{1}{u_n}+\frac{1}{u_{n+1}})$条件收敛，并求和。
\par\textbf{思路} 每一个内部倒数在有限和中出现两次，符号相反。
\par\textbf{详细解答} 直接展开给
\[
S_N=\dfrac1{u_1}+\dfrac{(-1)^{N-1}}{u_{N+1}}.
\]
由$\frac{u_n}{n}\to1$得$\frac{1}{u_{N+1}}\to0$，故$S_N\to\frac{1}{u_1}$。
又$u_n>0$最终成立，且
\[
n\left|\dfrac1{u_n}+\dfrac1{u_{n+1}}\right|\longrightarrow2.
\]
绝对值级数与调和级数同敛散，故发散；原级数条件收敛。
\par\textbf{方法点评} 这里没有假设$u_n$或$\frac{1}{u_n}+\frac{1}{u_{n+1}}$单调，有限消去已经给出普通和。
\par\textbf{疑难点} 不能只凭“一般项等价于$\frac{2(-1)^{n-1}}{n}$”使用极限比较法；极限比较用于非负项。
\end{example}

\begin{example}{竞赛风格｜非线性递推产生的精确望远镜和}{难度：$\star\star\star$}\label{cse:ex:nonlinear}
\textbf{题目} $a_1=c>1$，$a_{n+1}=a_n^2-a_n+1$。求$\sum_{n\ge1}\frac{1}{a_n}$及余项。
\par\textbf{思路} 递推两边减1后可以因式分解。
\par\textbf{详细解答} $a_{n+1}-a_n=(a_n-1)^2>0$，且$a_n\ge c>1$。
若$a_n$有有限极限$L\ge c$，则$L=L^2-L+1$，得$L=1$，矛盾，故$a_n\to\infty$。
由$a_{n+1}-1=a_n(a_n-1)$，
\[
\dfrac1{a_n}=\dfrac1{a_n-1}-\dfrac1{a_{n+1}-1},
\qquad
\sum_{n=1}^N\dfrac1{a_n}
=\dfrac1{c-1}-\dfrac1{a_{N+1}-1}.
\]
因此和为$\frac{1}{c-1}$，尾和精确为$\frac{1}{a_{N+1}-1}$。
\par\textbf{方法点评} 非线性迭代可能比线性迭代更容易形成倒数望远镜和。
\par\textbf{疑难点} $c=1$是恒定解，原级数发散，不能把它代入和公式。
\end{example}

\subsection{有序级数与隐式比较量}\label{cse:implicit}
\begin{method}{从符号或隐式方程构造正项比较}
\textbf{步骤} 已知两级数收敛且项有序时，研究差级数；隐式关系中先证明变量趋0，再用确定阶数的余项决定比较对象。
\par\textbf{疑难点} 未证明变量趋0时不能展开；任意项的大小关系本身不能保证绝对收敛传递。
\end{method}
\begin{example}{考研风格｜有序且收敛时的绝对收敛等价}{难度：$\star\star$}\label{cse:ex:ordered}
\textbf{题目} 实数项满足$a_n\le b_n$，且$\sum a_n,\sum b_n$均收敛。证明一个绝对收敛当且仅当另一个绝对收敛。
\par\textbf{思路} 差项非负，使通常收敛自动变成绝对收敛。
\par\textbf{详细解答} 令$d_n=b_n-a_n\ge0$。有限部分和相减并取极限，知$\sum d_n$收敛。
由$|b_n|\le|a_n|+d_n$，若$\sum|a_n|<\infty$则$\sum|b_n|<\infty$；由$|a_n|\le|b_n|+d_n$反向同理。
\par\textbf{方法点评} 有序关系并非直接比较绝对值，而是先产生一个可和的非负差。
\par\textbf{疑难点} 两级数都收敛是必要的前提。仅取$a_n=0,b_n=\frac{1}{n}$即说明删去该前提后结论失败。
\end{example}

\begin{example}{竞赛风格｜指数隐式关系中消失的一阶项}{难度：$\star\star\star$}\label{cse:ex:exponential}
\textbf{题目} $a_n,b_n>0$，$\e^{a_n}=a_n+\e^{b_n}$，且$\sum a_n<\infty$。证明$\sum \frac{b_n}{a_n}<\infty$。
\par\textbf{思路} 把$b_n$显式表示为一个对数，比较$a_n^2$阶的主项。
\par\textbf{详细解答} $a_n\to0$，并且
$b_n=\ln(\e^{a_n}-a_n)$。在$t\to0$时，
\[
\e^t-t=1+\dfrac12t^2+\dfrac16t^3+O(t^4),\qquad
\ln(\e^t-t)=\dfrac12t^2+O(t^3).
\]
故$\frac{b_n}{a_n^2}\to\frac{1}{2}$。例如充分大的$n$有$0<\frac{b_n}{a_n}\le a_n$，与给定收敛正项级数比较即可。
\par\textbf{方法点评} 方程消去了一阶项，把$b_n$压到$a_n^2$尺度。
\par\textbf{疑难点} 只用$b_n\to0$不能得到结论，需要证明$\frac{b_n}{a_n}$的具体阶数。
\end{example}

\begin{example}{竞赛风格｜余弦隐式关系不必求逆函数}{难度：$\star\star\star$}\label{cse:ex:cosine}
\textbf{题目} $0<a_n,b_n<\frac{\pi}{2}$，$\cos a_n-a_n=\cos b_n$，且$\sum b_n<\infty$。证明$\sum \frac{a_n}{b_n}<\infty$。
\par\textbf{思路} 保留方程整体，分别除以变量的自然主阶。
\par\textbf{详细解答} $\cos a_n>\cos b_n$且余弦在$(0,\frac{\pi}{2})$递减，故$0<a_n<b_n\to0$。
方程整理为$a_n+1-\cos a_n=1-\cos b_n$，所以
\[
\dfrac{a_n}{b_n^2}
=\dfrac{\frac{1-\cos b_n}{b_n^2}}{1+\frac{1-\cos a_n}{a_n}}
\longrightarrow\dfrac12.
\]
于是$\frac{a_n}{b_n}\le b_n$最终成立，比较判别得所求级数收敛。
\par\textbf{方法点评} 比值分解绕开$\arccos$的边界导数，控制更直接。
\par\textbf{疑难点} $\cos$在0处导数为0，因此不能按普通可微反函数的一阶公式处理。
\end{example}

\subsection{根列与振荡相位定义的级数}\label{cse:rootphase}
\begin{method}{将抽象一般项转成可控尺度}
\textbf{步骤} 根列先用相邻有理点夹住；振荡相位先减去整数倍$\pi$，再判断小相位的正性、单调性与阶数。
\par\textbf{疑难点} 条件收敛必须同时给普通收敛证明和绝对值发散证明。
\end{method}
\begin{example}{考研风格｜方程根定义的交错级数}{难度：$\star\star\star$}\label{cse:ex:rootseries}
\textbf{题目} $a_n$为$x^n+nx-1=0$的唯一正根。证明$\sum(-1)^{n-1}a_n$条件收敛，并比较$a_n$与$\frac{1}{n}$的误差。
\par\textbf{思路} 在$\frac{1}{n+1}$与$\frac{1}{n}$处计算函数符号。
\par\textbf{详细解答} 左端对$x>0$严格递增，在0处为$-1$，趋于无穷，故唯一正根存在。
对$n\ge2$，
\[
F_n(\dfrac{1}{n+1})=(n+1)^{-n}-(n+1)^{-1}<0,\qquad
F_n(\dfrac{1}{n})=n^{-n}>0.
\]
于是$\frac{1}{n+1}<a_n<\frac{1}{n}$。相邻根的区间分离给$a_{n+1}<a_n$，且$a_1=\frac{1}{2}>a_2$。
因此$a_n\downarrow0$，交错级数收敛；而$\sum a_n$由下界调和比较发散。
方程还给$1-na_n=a_n^n\le n^{-n}$，故
\[
n^{n+1}(\dfrac{1}{n}-a_n)=(na_n)^n=(1-a_n^n)^n\longrightarrow1,
\]
最后一步由$0\le1-(1-a_n^n)^n\le n a_n^n\le n^{1-n}\to0$。
\par\textbf{方法点评} 互不交叠的根区间能同时证明单调性和主阶，比连续化$n$求导更简便。
\par\textbf{疑难点} 各$n$分别存在唯一根，不自动保证根列单调，仍须比较。
\end{example}

\begin{example}{竞赛风格｜从整数倍相位中剥离交错因子}{难度：$\star\star\star$}\label{cse:ex:phase}
\textbf{题目} 对实参数$a$判定$\sum_{n\ge1}\sin(\pi\sqrt{n^2+a^2})$的敛散性。
\par\textbf{思路} 大相位是$n\pi$加一个小正修正量。
\par\textbf{详细解答} $a=0$时每项为0。若$a\ne0$，有
\[
\delta_n=\pi(\sqrt{n^2+a^2}-n)
=\dfrac{\pi a^2}{\sqrt{n^2+a^2}+n}\downarrow0,\qquad
\sin(\pi\sqrt{n^2+a^2})=(-1)^n\sin\delta_n.
\]
充分大的$n$有$0<\delta_n<\frac{\pi}{2}$，此后$\sin\delta_n$递减趋0，Leibniz判别给普通收敛。
又$n\sin\delta_n\to\frac{\pi a^2}{2}>0$，绝对值级数与调和级数比较发散，故条件收敛。
\par\textbf{方法点评} 有理化先得到小相位的精确表达，再讨论振荡，不会将整个正弦错当成小角。
\par\textbf{疑难点} 最初若干项的$\sin\delta_n$可能不单调，只需从足够大的指标起成立即可。
\end{example}

\subsection{离散分部积分保留边界项}\label{cse:weighted}
\begin{method}{加权差分还原部分和}
\textbf{条件} 已知$\sum n(a_n-a_{n-1})$和边界项$na_n$的行为。
\par\textbf{步骤} 在有限和层面移位，准确记录首项、末项；最后利用已有收敛条件。
\par\textbf{疑难点} 权重$n$不能在无穷和中随意移位，边界项可能决定全部结论。
\end{method}
\begin{example}{竞赛风格｜差分级数与边界项的联立条件}{难度：$\star\star\star$}\label{cse:ex:boundary}
\textbf{题目} 若$na_n\to A\in\R$，且$\sum_{n=2}^{\infty}n(a_n-a_{n-1})=T$收敛，证明$\sum a_n$收敛，并确定$A$。
\par\textbf{思路} 有限加权差分恰好等于边界项减去原部分和。
\par\textbf{详细解答} 对$N\ge2$，
\[
T_N=\sum_{n=2}^{N}n(a_n-a_{n-1})
=Na_N-2a_1-\sum_{n=2}^{N-1}a_n
=Na_N-a_1-\sum_{n=1}^{N-1}a_n.
\]
所以$\sum_{n=1}^{N-1}a_n=Na_N-a_1-T_N\to A-a_1-T$，原级数收敛。
若$A\ne0$，则$a_n\sim \frac{A}{n}$，最终同号且其绝对值不小于$\frac{|A|}{2n}$，与原级数收敛矛盾。因此$A=0$，和为$-a_1-T$。
\par\textbf{方法点评} 两个给定条件联用不仅推出级数收敛，还强迫边界极限为0。
\par\textbf{疑难点} 直接删去$Na_N$会在尚未证明$A=0$时循环论证。
\end{example}

\section*{章末分层习题}\addcontentsline{toc}{section}{章末分层习题}
\refstepcounter{exercise}\label{q:9:1}\par\noindent\textbf{习题\theexercise\quad 基础题} 判断$\sum\frac{1}{n^2}$。\quad\hyperref[a:9:1]{[答案]}\par
\refstepcounter{exercise}\label{q:9:2}\par\noindent\textbf{习题\theexercise\quad 基础题} 判断$\sum \frac{n}{n+1}$。\quad\hyperref[a:9:2]{[答案]}\par
\refstepcounter{exercise}\label{q:9:3}\par\noindent\textbf{习题\theexercise\quad 基础题} 求$\sum_{n=0}^{\infty}(\frac{1}{3})^n$。\quad\hyperref[a:9:3]{[答案]}\par
\refstepcounter{exercise}\label{q:9:4}\par\noindent\textbf{习题\theexercise\quad 考研风格} 判断$\sum\frac{(-1)^{n-1}}{n^2}$是绝对还是条件收敛。\quad\hyperref[a:9:4]{[答案]}\par
\refstepcounter{exercise}\label{q:9:5}\par\noindent\textbf{习题\theexercise\quad 考研风格} 求$\sum_{n\ge1}\frac{x^n}{n^2}$的收敛域。\quad\hyperref[a:9:5]{[答案]}\par
\refstepcounter{exercise}\label{q:9:6}\par\noindent\textbf{习题\theexercise\quad 考研风格} 求$\sum_{n\ge1}\frac{n}{3^n}$。\quad\hyperref[a:9:6]{[答案]}\par
\refstepcounter{exercise}\label{q:9:7}\par\noindent\textbf{习题\theexercise\quad 竞赛拓展} 证明$\sum \frac{x^n}{n^2}$在$[-1,1]$上一致收敛。\quad\hyperref[a:9:7]{[答案]}\par
\refstepcounter{exercise}\label{q:9:8}\par\noindent\textbf{习题\theexercise\quad 竞赛拓展} 求$\sum_{n\ge1}\frac{1}{(2n-1)^2}$。\quad\hyperref[a:9:8]{[答案]}\par
\refstepcounter{exercise}\label{q:9:9}\par\noindent\textbf{习题\theexercise\quad 竞赛拓展} 求$\sum_{n\ge1}\frac{1}{n(n+2)}$。\quad\hyperref[a:9:9]{[答案]}\par
\section*{章末方法速查}\addcontentsline{toc}{section}{章末方法速查}
\par\noindent\begin{tabularx}{\textwidth}{p{23mm}XXX}\toprule 方法&适用条件&典型信号&易错点\\\midrule
比较积分&正项&幂与对数&方向用反\\Leibniz&幅度单调趋零&交错符号&只看交错\\幂级数运算&半径内部&系数$n,\frac{1}{n}$&端点未查\\Fourier&周期与收敛条件&三角基与选点&跳点未取平均\\
\bottomrule\end{tabularx}
